% =====================================================================
% The 15:50 Candle -- full working paper.
% It serves as the Internet Appendix to the shorter journal version in
% ../spx_candle_1550_journal/. The bibliography is inline (no bibtex
% pass) and the figures are drawn from the CSV files in figures/.
% =====================================================================
\documentclass[12pt,letterpaper]{article}

\usepackage[utf8]{inputenc}
\usepackage[T1]{fontenc}
\usepackage{amsmath,amsthm}
\usepackage{lmodern}
\usepackage[margin=1in]{geometry}
\usepackage{setspace}
\usepackage{booktabs}
\usepackage{array}
\newcolumntype{L}[1]{>{\raggedright\arraybackslash}p{#1}}
\usepackage{microtype}
\usepackage{natbib}
\usepackage{multirow}
\usepackage{float}
\usepackage{xcolor}
\usepackage{etoolbox}
\usepackage{enumitem}
\usepackage[font=small,labelsep=colon,skip=4pt]{caption}
\captionsetup[table]{position=top}
\captionsetup[figure]{position=top}
\usepackage[compact]{titlesec}
\usepackage{titling}
\usepackage{placeins}
\usepackage{needspace}
\usepackage{fancyhdr}
\usepackage{pgfplots}
\pgfplotsset{compat=1.18}
\usepgfplotslibrary{groupplots}
\definecolor{linkblue}{HTML}{1C4E8C}
\usepackage[colorlinks=true,allcolors=linkblue,hyperfootnotes=false]{hyperref}
\hypersetup{pdftitle={The 15:50 Candle: A Pre-Registered Study of an End-of-Day Signal in 0DTE SPX Options}, pdfauthor={Kausthubh Veldanda}}

% ---- page style: page numbers only ----
\pagestyle{plain}

% compact title block: title near the top margin, author, date and abstract close below
\setlength{\droptitle}{-5.5em}
\pretitle{\begin{center}\Large}
\posttitle{\par\end{center}\vspace{-1.3em}}
\preauthor{\begin{center}\normalsize}
\postauthor{\par\end{center}\vspace{-1.5em}}
\predate{\begin{center}\small}
\postdate{\par\end{center}\vspace{-2.4em}}

\setlist{itemsep=1pt,topsep=3pt}

\setcounter{topnumber}{3}
\setcounter{bottomnumber}{2}
\setcounter{totalnumber}{5}
\renewcommand{\topfraction}{0.9}
\renewcommand{\bottomfraction}{0.8}
\renewcommand{\textfraction}{0.07}
\renewcommand{\floatpagefraction}{0.8}
\setlength{\textfloatsep}{10pt plus 4pt minus 3pt}
\setlength{\floatsep}{8pt plus 2pt minus 2pt}
% pages that hold only tables or figures: stack them from the top, no spreading
\makeatletter
\setlength{\@fptop}{0pt}
\setlength{\@fpsep}{16pt}
\setlength{\@fpbot}{0pt plus 1fil}
\makeatother
\setlength{\parskip}{0pt plus 1pt}
\flushbottom
\setlength{\intextsep}{8pt plus 4pt minus 2pt}
% table body size; a wide table sets \tabfont to 9pt locally
\newcommand{\tabfont}{\footnotesize}
\AtBeginEnvironment{tabular}{\tabfont}
% tables and figures single-spaced
\AtBeginEnvironment{table}{\setstretch{1}}
\AtBeginEnvironment{figure}{\setstretch{1}\setlength{\parindent}{0pt}}
\setlength{\emergencystretch}{2em}

% the explanatory note of a table or figure, set below it in small type
\newcommand{\tabnotes}[1]{\par\vspace{5pt}\parbox{\linewidth}{\footnotesize\setlength{\parindent}{0pt}\textit{Note:} #1}\par}

% ---- figures: drawn by pgfplots from CSV files made by figures/make_figure_data.py ----
\definecolor{series1}{HTML}{2A78D6}
\definecolor{series2}{HTML}{EB6834}
\pgfplotsset{
  table/col sep=comma,
  every axis/.append style={font=\fontsize{9}{11}\selectfont},
  paper/.style={
    axis lines*=left, axis line style={black!60}, tick style={black!60},
    tick align=outside, major tick length=2pt,
    ymajorgrids, grid style={black!10},
    label style={font=\fontsize{10}{12}\selectfont},
    legend style={draw=none, fill=none, font=\fontsize{9}{11}\selectfont, legend cell align=left},
    /pgf/number format/1000 sep={,}, scaled y ticks=false,
    y tick label style={/pgf/number format/fixed},
  },
  zeroline/.style={extra y ticks={0}, extra y tick labels={},
    extra y tick style={grid=major, major grid style={black!55, thin}, tick style={draw=none}}},
}

\newcommand{\ITM}[1]{\textsc{itm}#1}
\newcommand{\OTM}[1]{\textsc{otm}#1}
\newcommand{\ATM}{\textsc{atm}}
\newcommand{\dollar}{\$}

\begin{document}
\title{\textbf{The 15:50 Candle: A Pre-Registered Study\\of an End-of-Day Signal in 0DTE SPX Options}\thanks{This version: 12 November 2026.
Every confirmatory test in this paper was pre-registered: each plan was hashed and
time-stamped before the data it governs was read, and the tests that failed are
reported with those that passed. This is the full working-paper version of the
study and reports every study in full; it serves as the Internet Appendix to the
shorter journal version of this paper. A statement on data availability is at the end of the paper. I thank Tiger
Chen, whose observation started this study. This paper is research, not
investment advice.}}
\author{Kausthubh Veldanda\thanks{University of California, Berkeley. Email: \href{mailto:kveldanda@berkeley.edu}{kveldanda@berkeley.edu}}}
\date{12 November 2026}
\maketitle

\begin{abstract}
\noindent I test whether a one-minute candle near the close predicts the move of the
S\&P~500 index (SPX) into the close, and whether the prediction is worth
anything after the cost of trading same-day-expiry (0DTE) index options. Every confirmatory test was
pre-registered. A scan of 41 candle features on 2013--2020 selected a 15:50~ET
candle whose body spans at least 60\% of its range. On the 2021--2026 holdout,
traded in its own direction, it was right on 57.5\% of 644 trades. At real OPRA
quotes, an SPXW option 10 points in the money, bought at the ask and held to
settlement, earned \$67 a trade (Newey--West $t = 2.40$), about two-thirds of what
a pricing model implied. The profit depends on the bid-ask spread, and the same trades lost money before
2021, when spreads were about five times wider. Early exits
(including 26 pre-registered stop and take-profit rules at real bids), limit
orders, momentum and dealer-gamma filters and smaller contracts did not improve
the trade. The result also depends on where the candle is read. Read from the SPX index
candle instead of the CFD candle of the backtests, the rule made \$39 a trade
($t = 1.53$) and missed its criterion. Read from the SPY candle it made \$72
($t = 2.35$) and met a retrospective one. A final search over 7{,}128 settings of
candle, cut, direction and strike again chose the 15:50 candle, with a luck
diagnostic of $p = 0.051$. In a pre-registered test using the exchanges' imbalance messages, the candle pointed the same way as the closing-auction imbalance
published at 15:50 on 73\% of signal days but was not more accurate on those days
(56\% against 62\%). I read the evidence as a small edge that is sensitive to execution and not yet
established.

\medskip
\noindent \textit{JEL:} G13, G14\\
\noindent \textit{Keywords:} 0DTE options; S\&P~500 index options; closing
auction; intraday momentum; bid-ask spread; pre-registration.
\end{abstract}



% =====================================================================
\section{Introduction}
\label{sec:intro}

Options on the S\&P~500 index that expire the same day (0DTE) are now listed
every weekday, and much of their volume concentrates near the close. The
closing auction also becomes visible in the last ten minutes of the session.
From 15:50~ET the New York Stock Exchange publishes the closing-auction imbalance of the stocks
it lists, and Nasdaq disseminates an early imbalance indicator for its own
stocks. An index option bought at 15:51~ET settles in cash on the official close
nine minutes later, so its outcome depends on little more than the direction and
size of the last move of the day. The close is therefore a reasonable place to look for a short-horizon
directional signal, and a 0DTE option is the simplest way to trade one.

I came to this question through a friend, Tiger Chen, who had noticed by
chance, on one-minute charts of the index, that the S\&P~500 often makes a
sharp move in one direction about ten minutes before the close, and who showed
me 27 screenshots of it. This paper began as an attempt to check whether that
observation is real and whether trading it with same-day options makes money.

A backtest of such a trade can mislead in several ways. If enough candle rules
are scored on enough history, one of them will look good by chance, and its
edge will shrink out of sample \citep{harvey2016}. A backtest that enters at
the close of the signal candle uses information that was not yet available. And
a rule that is right 54\% of the time can still lose money once the bid-ask
spread is paid. I therefore fixed what would be tested before reading the data,
and the confirmation prices the trade at quotes that were available at the
time.

Every confirmatory test in the paper was pre-registered \citep{nosek2018}.
Before reading the data for a study I wrote down what would be computed, on
which data, and what would count as success, and the plan was hashed and
time-stamped. The tests that failed are reported with those that passed. The
studies were done in sequence, each building on the last. A scan of 41 candle features finds a signal on 2013--2020
and tests it on 2021--2026 (Section~\ref{sec:scan}). A calibrated pricing model
then asks which strike best expresses the signal and whether a stop-loss or
take-profit beats holding to settlement (Section~\ref{sec:model}). Real OPRA
quotes replace the model to tune and confirm the trade at the prices actually
offered (Section~\ref{sec:realprice}). Trade bars describe how the order might
have been filled, and a pre-registered study tests real-time exit rules at real
bids (Section~\ref{sec:execution}). Three further studies re-tune the rule on
recent data (Section~\ref{sec:refine}), test published intraday momentum and a
dealer-gamma channel (Section~\ref{sec:momentum}), and trade the signal on
smaller contracts (Section~\ref{sec:mini}). An independent re-validation then
checks the whole chain (Section~\ref{sec:audit}). Because the backtests read the
signal from a contract for difference (CFD) that tracks the index, two further
pre-registered studies read the signal from the SPX index candle itself
(Section~\ref{sec:realindex}) and from the SPY candle
(Section~\ref{sec:spycandle}), and another searches 7{,}128 nearby settings,
including other closing candles, for a better rule (Section~\ref{sec:search}). A
last study tests an explanation of the signal with the closing-auction imbalances
that the exchanges publish at 15:50 (Section~\ref{sec:imbalance}).

The scan produced one usable signal, the 15:50 one-minute candle with a
\emph{strong body}, traded in its own direction. On 2021--2026 it was right on
57.5\% of 644 trades. It is a stricter version of an effect I had already seen
in earlier exploratory work, and it has weakened since 2023.

At real option prices the edge is small. Deep in-the-money options did best, in
the model and at real prices. An option 10 points in the money, bought at the
quoted ask and held to settlement, earned $+\$67.46$ per contract on 2021--2026
($t = 2.40$), about two-thirds of the model's estimate. Whether the trade makes
money depends mostly on the bid-ask spread. The same trades lost money at the
ask before 2021, and a limit order at the mid does not recover the spread,
because it misses the best days. No early exit improved on holding to
settlement, including 26 stop and take-profit rules tested at real bids.
Re-tuned thresholds, momentum filters and cheaper contracts did not improve on
the original rule either.

The result also depends on where the candle is read. An independent
re-validation reproduced the results above, placed the start of the effect near a
60\% body, and found that a larger body beyond that threshold does not mean a
larger gain. It also showed that the CFD candle from which the signals were
computed reads weaker than the SPX index candle. Read from the index candle,
under a plan locked before any index candle was joined with a price, the rule
made $+\$39.44$ a trade ($t = 1.53$) and did not meet its criterion. The gain was
confined to days on which both candles showed a strong body, so the confirmation
above is a statement about the CFD candle. Read from the SPY candle, which can be
seen live, the rule made $+\$71.74$ a trade ($t = 2.35$) and met its criterion.
SPY was the third source evaluated on the same outcomes, however, so this
result is retrospective and is not an independent confirmation, and the gain
was
again confined to sessions on which all three candles showed a strong body.

A later pre-registered search over 7{,}128 settings of candle source, minute, cut, direction and strike chose
the 15:50 candle again. Its choice was not clearly beyond luck ($p = 0.051$), and
because 15:50 settings had already been examined on the later years, it could
not be tested afresh. Finally, the candle pointed the same way as the
closing-auction imbalance published at 15:50 on 73\% of signal days, but in a
pre-registered test it was not more accurate on those days (56\% against 62\%).
The explanation that traders react to the published imbalance was therefore not
confirmed, although the test could detect only a large difference.

Several limits apply throughout, and none of this is a recommendation to trade.
The holdout was not blind, because a related effect was already known. The
strike and exit results rest on a model calibrated on seven days. The
real-price confirmation reuses index outcomes that I had already seen, and the
execution results are descriptive proxies. The confirmed rule reads its candle
from a CFD and not from the index, and the SPY candle's pass is retrospective. Section~\ref{sec:limits} discusses these limitations.

\paragraph{Related literature.} The closest work is on market intraday momentum. \citet{ghlz2018} find that the first
half-hour return predicts the last half hour of the S\&P~500, and
\citet{bdlm2021} find the same pattern across futures markets and link it to the
hedging of option dealers who are short gamma. The 15:50 candle lies inside that
last half hour, and I test both the momentum and the gamma channel directly. What I add is a
single minute within that window, tested out of sample under a pre-registered
plan. The holdout design and the luck diagnostics follow the work on data
mining and multiple testing in return predictability \citep{harvey2016}.
Pre-registration \citep{nosek2018} is still rare in empirical finance, and the
paper applies a plan, hash and ledger protocol to a sequence of trading studies
and reports the tests that failed.

\paragraph{Outline.} Section~\ref{sec:data} describes the data and
Section~\ref{sec:design} the research design. Section~\ref{sec:scan} tests the
signal, and Section~\ref{sec:model} studies strike choice and early exits under a
pricing model. Sections~\ref{sec:realprice} and~\ref{sec:execution} price the
trade at real quotes and examine how it would have been executed.
Sections~\ref{sec:refine} to~\ref{sec:mini} report the studies that tried to
improve the rule, and Section~\ref{sec:audit} the re-validation.
Sections~\ref{sec:realindex} and~\ref{sec:spycandle} read the candle from the
index and from SPY, Section~\ref{sec:search} searches the other closing candles,
and Section~\ref{sec:imbalance} tests the closing-auction explanation.
Section~\ref{sec:forward} reports the forward evidence,
Section~\ref{sec:discussion} discusses what the studies imply together, and
Section~\ref{sec:conclusions} concludes.

% =====================================================================
\section{Data}
\label{sec:data}

The study combines one-minute bars for the index and for SPY, real option
quotes and trades, the exchanges' closing-auction imbalance messages and a
broker's own one-minute bars. Table~\ref{tab:data} lists the sources.

\begin{table}[htbp]
\centering
\caption{Data sources.}
\label{tab:data}
\begin{tabular}{@{}L{0.17\linewidth}L{0.41\linewidth}L{0.13\linewidth}L{0.21\linewidth}@{}}
\toprule
Source & Content & Period & Used for \\
\midrule
Dukascopy & US500 CFD one-minute bars, regular trading hours & 2013--2026 & Signals; index outcomes \\
Databento OPRA & SPXW best bid and offer at one-minute marks, 15:00--16:01~ET, the 25 strikes nearest the money, calls and puts & 2013--2026 & Option prices (Sections~\ref{sec:realprice}--\ref{sec:refine}) \\
Databento OPRA & SPXW one-minute trade bars for the same contracts & 2013--2026 & Execution (Section~\ref{sec:execution}) \\
Databento OPRA & Daily open interest, all SPX and SPXW options & 2013--2026 & Gamma proxy (Section~\ref{sec:momentum}) \\
Databento OPRA & XSP and SPY quotes; XSP trade bars & 2021--2026 & Smaller contracts (Section~\ref{sec:mini}) \\
Cboe & Official SPX daily close & 2013--2026 & Option settlement \\
Robinhood & SPX and SPXW one-minute bars & 2026 & Model calibration; forward test \\
CNBC & SPX index one-minute bars (archive) & 2021--2026 & Index candle (Section~\ref{sec:realindex}) \\
CNBC & SPY one-minute bars (archive; traded prices, not dividend-adjusted) & 2021--2026 & SPY candle (Section~\ref{sec:spycandle}) \\
Databento & NYSE and Nasdaq closing-auction imbalance messages, 15:50--16:00~ET on signal days, for SPY's 50 largest holdings (SEC N-PORT reports) & 2021--2026 & Imbalance test (Section~\ref{sec:imbalance}) \\
\bottomrule
\end{tabular}
\tabnotes{The index sample runs from 2013-02-28 to 2026-10-20 and the OPRA SPXW
sample from 2013-04-01 to 2026-10-20; the XSP and SPY sample starts on
2021-01-01. Robinhood data begin on 2026-09-21. The CFD series was extended to 2026-11-09
for the comparison in Section~\ref{sec:audit}.}
\end{table}

Signals and index outcomes come from Dukascopy's US500 CFD one-minute bars during
regular trading hours, restricted to a validated set of full sessions. Discovery
uses 2013--2020 (1{,}861 sessions) and the holdout 2021--2026 (1{,}446
sessions). The CFD stands in for SPX, and index outcomes are measured to the
close of the 15:59 bar rather than to the official settlement.
Section~\ref{sec:realindex} reads the signal candle instead from an archive of
SPX index one-minute bars, and Section~\ref{sec:spycandle} from an archive of SPY
one-minute bars. On the 22 days where both are available, the CFD and SPX moves
from 15:51 to the close had the same sign on 86\% of days.

Real option prices come from Databento's OPRA feed: the
consolidated best bid and offer at one-minute marks for SPXW contracts expiring
that day, from 15:00 to 16:01~ET, for the 25 strikes within 60 points of the
at-the-money strike at 15:51. A quote stamped HH:MM:00 is the quote at the open
of that minute. One-minute trade bars for the same contracts support the
execution analysis, and daily open interest on every SPX and SPXW option
supports the dealer-gamma proxy. Of the 3{,}287 index days, 2{,}149 had an SPXW
expiry that day; the rest are excluded from every option study. Options settle at
intrinsic value on Cboe's official SPX close, which matched the broker's daily
close on all 362 days on which both exist.

Finally, Robinhood's SPX and SPXW bars serve two purposes. Seven
days of real one-minute option prices calibrate the pricing model of
Section~\ref{sec:model}, and since 21 October 2026 a read-only collector has
saved the day's SPX and SPXW one-minute bars after every close for the forward
test (Section~\ref{sec:forward}). The broker keeps option minute bars for only
about a day, so a day that is not collected promptly is lost.

% =====================================================================
\section{Research design}
\label{sec:design}



\subsection{The signal and the trade}
\label{sec:signal}

Let $o, h, l, c$ be the open, high, low and close of a one-minute bar. The bar
has a \emph{strong body} if $h > l$ and
\begin{equation}
|c - o| \ge 0.6\,(h - l),
\end{equation}
and its direction is the sign of $c - o$. The signal studied throughout is a
strong body on the 15:50 bar. Because that bar is complete only at 15:50:59, the
trade enters at the open of the 15:51 bar and, in the index tests, is scored on
the move to the close of the 15:59 bar.

In the option studies, let $S$ be the index level at entry and
$K_{\textsc{atm}} = 5\lfloor S/5 + 1/2\rfloor$ the nearest listed strike. An
offset of $j$ points, positive meaning in the money in the trade's direction,
buys the call struck at $K_{\textsc{atm}} - j$ after an up signal and the put
struck at $K_{\textsc{atm}} + j$ after a down signal. I consider nine offsets,
from 20 points in the money (\ITM{20}) to 20 points out of the money
(\OTM{20}) in steps of five. One contract (multiplier 100) is bought per signal
and held to cash settlement unless stated otherwise, and each side pays a fee of
\$0.09 per contract, with \$0.74 as a sensitivity. I refer to the 15:50 strong-body signal traded at \ITM{10} as the \emph{current
rule}.

\subsection{Pre-registration}
\label{sec:prereg}

Each confirmatory study was specified in a written plan before the data it
governs was read. The plan's SHA-256 hash and a UTC timestamp were appended to a
ledger, and each study's code recomputes the hash at start-up and refuses to
read real data unless it matches a ledger entry. Each plan lists what was already
known when it was written, and any departure from a plan is reported as a
labelled deviation; no result was re-run under a changed rule.
Table~\ref{tab:studies} lists the studies, their primary criteria and their
outcomes. Analyses that were not pre-registered are labelled \emph{descriptive} and were
not used to decide any test.

\begin{table}[htbp]
\centering
\caption{The pre-registered studies.}
\label{tab:studies}
\begin{tabular}{@{}L{0.16\linewidth}L{0.07\linewidth}L{0.11\linewidth}L{0.33\linewidth}L{0.22\linewidth}@{}}
\toprule
Study & Section & Registered & Primary criterion & Outcome \\
\midrule
Candle scan & \ref{sec:scan} & 2026-10-21 & Holm-corrected test of each frozen candidate on 2021--2026, and an accuracy bar & One of three candidates passes \\
Strike ladder & \ref{sec:strikes} & 2026-11-08 & Nine strikes ranked under a calibrated model & Deep in the money most robust \\
Early exits & \ref{sec:exits} & 2026-11-09 & A rule beats holding with $t > 2.64$, robust across model variants & No rule passes \\
Real option prices & \ref{sec:realprice} & 2026-11-09 & Switch only if a new setting leads by $t \ge 2.0$; confirm with $t \ge 2.0$ & Current rule kept and confirmed \\
Re-tuning & \ref{sec:refine} & 2026-11-09 & Luck diagnostic $p \le 0.05$, then $t \ge 2.0$ on 2024--2026 & Pick fails the luck diagnostic \\
Intraday momentum & \ref{sec:momentum} & 2026-11-09 & Slope $t \ge 2.0$; filters $p \le 0.0125$; gamma $p \le 0.025$ & Nothing supported \\
XSP and SPY & \ref{sec:mini} & 2026-11-09 & Positive mean with $t \ge 2.0$ on 2021--2026 & Neither confirmed \\
Index candle & \ref{sec:realindex} & 2026-11-10 & Rule read from the index candle: positive mean per day with $t \ge 2.0$ & Criterion not met \\
SPY candle & \ref{sec:spycandle} & 2026-11-10 & Rule read from the SPY candle: positive mean per day with $t \ge 2.0$ & Retrospective criterion met \\
Candle search & \ref{sec:search} & 2026-11-10 & Best of 7{,}128 settings on 2021--2023, tested on 2024--2026 only if not at 15:50 and luck diagnostic $p \le 0.05$ & Pick at 15:50; no test \\
Closing imbalance & \ref{sec:imbalance} & 2026-11-10 & Hit rate higher on days the candle agrees with the 15:50 imbalance: $d > 0$ with $t \ge 2.0$ & Not confirmed \\
Real-time exits & \ref{sec:rtexit} & 2026-11-10 & An exit rule beats holding in both halves of 2021--2023 and has luck diagnostic $p \le 0.05$; then $t \ge 2.0$ on 2024--2026 & Stopped at the coverage gate \\
Real-time exits, revised & \ref{sec:rtexit} & 2026-11-11 & The same, for the 26 rules that read only the option's bid & No eligible rule; no test \\
Forward options log & \ref{sec:forward} & 2026-10-30 & Running record of real option outcomes & Four priced trades so far \\
\bottomrule
\end{tabular}
\tabnotes{Registration dates are UTC. The limit-order and fixed-minute sale analyses of
Section~\ref{sec:execution} and the re-validation of Section~\ref{sec:audit}
were not pre-registered.}
\end{table}

\subsection{Statistical inference}
\label{sec:inference}

I report mean P\&L per trade and per day, where the daily series takes the
trade's P\&L on a signal day and zero otherwise. Unless stated otherwise,
$t$-statistics use Newey--West standard errors with five lags on the daily
series \citep{neweywest1987}. The scan tests its candidates one-sided with a
Holm correction \citep{holm1979}, and hit-rate intervals are Wilson intervals
\citep{wilson1927}.

Two studies choose one setting from a large grid, and for them I report a
\emph{luck diagnostic}. The signal series is shifted circularly against the
outcome days, by every shift from 21 days to $n - 21$ in steps of three. A shift
keeps the signal's own structure but breaks its link to the outcomes. The whole
search is repeated for each shift and its best score kept. With $S$ shifts, of
which $k$ produce a best score at least as good as the real one,
$p = (1 + k)/(1 + S)$. The diagnostic asks whether the real winner beats the
winner that unrelated signals would find. It is approximate, because expiry
frequency, volatility and pricing change over the sample.

\subsection{Option prices}
\label{sec:pricing}

The first option studies had no historical option prices and priced each trade
with a Bachelier model \citep{bachelier1900}:
\begin{equation}
C = (S-K)\,\Phi(d) + \sigma\,\phi(d), \qquad P = C - (S-K), \qquad d = \frac{S-K}{\sigma},
\end{equation}
where $\sigma$ is the standard deviation, in index points, of the move from entry
to the close. Each day's $\sigma$ is proportional to that day's one-minute
realised volatility, scaled to the nine minutes to the close. The constant, and a
skew factor for each strike, were calibrated on seven days of real Robinhood
SPXW prices (Appendix~\ref{app:cal}). The model charges a \$0.10 half-spread on
entry. To compare strikes on equal terms, every modelled trade is rescaled to a
common index level of 7{,}709.58 before it is priced, so model results are in
\emph{level-normalised} dollars. The later studies replace the model with real
OPRA quotes, buying at the ask at the open of the entry minute; their results
are in dollars at the index level of the time.

% =====================================================================
\section{Does a closing candle predict the close?}
\label{sec:scan}



I defined 41 features, each of which outputs up, down or no
trade using only bars completed before the entry: the direction of every
one-minute candle from 15:30 to 15:50; five candle shapes (strong body, hammer,
shooting star, engulfing, big candle) on the 15:50 one-minute candle and on the
15:45--15:49 five-minute candle; five-minute directions; and four whole-day
features. Features with $|t| \ge 2.5$ on 2013--2020 qualified, up to eight. The
list of qualifying features was frozen and hashed before any 2021 bar was read.
On 2021--2026 each candidate faced a one-sided test with a Holm correction at
5\%. A candidate that passed was called \emph{accurate} only if the lower bound
of its 95\% Wilson interval exceeded the hit rate at which buying the
at-the-money option breaks even.

\begin{table}[htbp]
\centering
\caption{The candle scan: frozen candidates on the 2021--2026 holdout.}
\label{tab:scan}
\begin{minipage}{\linewidth}\centering
\textit{Panel A. The three candidates}\par\smallskip
\setlength{\tabcolsep}{4pt}
\begin{tabular}{lcccccc}
\toprule
Candidate & 2013--20 & 2021--26 (95\% CI) & trades & Holm $p$ & predicts & accurate \\
\midrule
\textbf{15:50 strong body} & 54.1\% & \textbf{57.5\%} (53.6--61.2) & 644 & $6.3\times10^{-5}$ & yes & \textbf{yes, just} \\
15:50 direction, any body & 52.2\% & 54.0\% (51.4--56.6) & 1{,}445 & 0.011 & yes & no \\
Day so far (09:30--15:49) & 52.0\% & 48.4\% (45.8--51.0) & 1{,}446 & 0.53 & no & no \\
\midrule
\emph{P3 reference (not tested)} & \emph{51.8\%} & \emph{53.6\%} (51.1--56.2) & \emph{1{,}442} & --- & --- & --- \\
\bottomrule
\end{tabular}
\end{minipage}

\bigskip
\begin{minipage}{\linewidth}\centering
\textit{Panel B. The strong-body candle by subperiod (descriptive)}\par\smallskip
\begin{tabular}{lcccc}
\toprule
Strong body & trades & hit rate (95\% CI) & $t$ & option P\&L / trade ($t$) \\
\midrule
2021--2022 & 216 & 63.0\% (56.3--69.1) & 3.56 & $+\dollar80.7$ (2.33) \\
2023--2026 & 428 & 54.7\% (49.9--59.3) & 2.27 & $+\dollar28.6$ (1.11) \\
\bottomrule
\end{tabular}
\end{minipage}
\tabnotes{Panel A gives hit rates in discovery (2013--2020) and on the holdout,
with 95\% Wilson intervals. The strong-body break-even is 49.5\% at a \$0.09 fee
(51.8\% size-neutral); with 644 trades, ``accurate'' requires an observed hit
rate of 53.4\% (55.7\% size-neutral). The P3 row is a reference and was not
tested. Panel B prices the at-the-money option with the scan's cost model. By
year from 2021 to 2026 the hit rate was 64.9, 61.0, 54.0, 53.8, 57.3 and
53.3\%.}
\end{table}

Panel A of Table~\ref{tab:scan} shows that three features qualified. Two of them, both on the 15:50 candle, predicted the 2021--2026 move, but only
the strong-body candle also cleared the accuracy bar. It was right on 57.5\% of
644 trades. Rescaled to today's index level, these trades moved an average of
$+1.82$ points in the signal's direction over nine minutes. The third candidate,
a whole-day feature, failed out of sample.

I ran two checks, one against chance and one against a generic intraday
pattern. When the 41 features were shifted out of line with the
2013--2020 outcomes and the scan was re-run 304 times, about 1\% of runs found as
many qualifying features as the real scan, and none found one as strong as the
best real feature. The same rules moved four hours earlier, entering at 11:51,
showed nothing in either period.

\label{sec:scan-known}%
The result is not entirely new. Both passing features sit on the 15:50 candle
and agree with the direction of the move from 15:50 to 15:51 (an effect I call
P3, found in the project's earlier exploratory work) on 98\% and 93\% of days.
I knew how P3 had behaved over 2013--2026 before running the scan, so the
holdout was not blind for these features. I therefore read the result as a
known effect passing a stricter test, not as a new discovery. The strong body
does matter, because on 15:50 candles without one, trading the direction earned
nothing.

The effect has also weakened since 2023. Panel B of Table~\ref{tab:scan} shows
that most of the edge comes from 2021--2022. In 2023--2026 the hit-rate interval includes
50\%, and the modelled option P\&L is not significant. At that effect size, a
test with 80\% power would need roughly 19 years of signals, so no single year
can settle the question.

% =====================================================================
\section{Strike choice and early exits under a pricing model}
\label{sec:model}

Before real option quotes were available, two pre-registered studies used the
calibrated model of Section~\ref{sec:pricing} to price the 644 holdout trades.
The first asks which strike best expresses the signal, and the second whether a
stop-loss or take-profit improves on holding to settlement. Both report a
headline calibration, variant (a), and alternatives that change the volatility,
the costs or the skew (Appendix~\ref{app:cal}). Under the model, strikes deep in the money were the most robust and no early
exit improved on holding.

\subsection{Which strike}
\label{sec:strikes}

\begin{table}[htbp]
\centering
\caption{The strike ladder under the calibrated model.}
\label{tab:ladder}
\begin{tabular}{lrrrrrr}
\toprule
Strike & model price & mean P\&L / contract & $t$ & profitable & mean return & median return \\
\midrule
\ITM{20} & \dollar2{,}036 & $+162.21$ & 3.89 & 56\% & $+8\%$ & $+6\%$ \\
\ITM{15} & \dollar1{,}580 & $+142.62$ & 3.50 & 55\% & $+9\%$ & $+6\%$ \\
\textbf{\ITM{10}} & \textbf{\dollar1{,}098} & $\mathbf{+166.08}$ & \textbf{4.25} & \textbf{54\%} & $\mathbf{+15\%}$ & $\mathbf{+8\%}$ \\
\ITM{5} & \dollar693 & $+160.97$ & 4.49 & 51\% & $+23\%$ & $+3\%$ \\
\ATM & \dollar367 & $+154.63$ & 4.96 & 42\% & $+45\%$ & $-50\%$ \\
\OTM{5} & \dollar205 & $+92.10$ & 3.59 & 27\% & $+53\%$ & $-100\%$ \\
\OTM{10} & \dollar204 & $-34.34$ & $-1.47$ & 14\% & $+7\%$ & $-100\%$ \\
\OTM{15} & \dollar120 & $-21.91$ & $-1.11$ & 9\% & $+13\%$ & $-100\%$ \\
\color{gray}\OTM{20} & \color{gray}\dollar1{,}683 & \color{gray}$-1{,}628.66$ & \color{gray}$-3.62$ & \color{gray}1\% & \color{gray}$-74\%$ & \color{gray}$-100\%$ \\
\bottomrule
\end{tabular}
\tabnotes{Variant (a), 644 modelled strong-body trades on 2021--2026, \$0.09
fee, level-normalised dollars per contract. A model, not observed P\&L. The
\OTM{20} row (grey) is an artefact of the skew rule; the \OTM{10} and \OTM{15}
rows inherit skew factors of 2.09 and 2.26 and are unreliable for the same
reason. Table~\ref{tab:variants} reports every variant.}
\end{table}

Table~\ref{tab:ladder} shows that strikes 10 to 20 points in the money were the
most robust. Their mean stayed between $+\$102$ and $+\$214$ per
contract in every credible variant, a majority of trades were profitable, and the
typical trade made money. Deep in-the-money options carry little time value, so they keep nearly all of a
small move. At-the-money options earn a higher return on premium, but from
fewer, larger winners. The median at-the-money trade lost
about half its premium, and the mean depends heavily on the assumed price of
volatility. The ranking held in both subperiods, although the level fell after
2022.

\label{sec:skewfail}%
The pre-registered skew factor, the median ratio of market to model price,
failed for far out-of-the-money strikes. On the calm calibration days the normal
model priced a 20-point out-of-the-money option at almost nothing, while the
market charged at least the minimum tick, so the factor came out at 54. Applied
to volatile days, it produced absurd prices. Because the rule was fixed in
advance, its output is reported unchanged and marked as unreliable. The real
prices of Section~\ref{sec:realprice} settle the question for those strikes.

\subsection{Early exits}
\label{sec:exits}

Six exit rules watch the index's one-minute closes from 15:51 to 15:58: stops at
3, 5 or 8 points against the trade (S3, S5, S8), take-profits at 5 or 10 points
in its favour (T5, T10), and a stop at 5 combined with a target at 10 (B5\_10). A
rule that fires at the close of a minute sells at the next minute's open, paying
the half-spread and a second fee. The baseline (H) holds to settlement. A rule
counts as better than holding only if, at \ITM{10} in variant (a), its mean
paired difference from holding is positive with $t > 2.64$ (two-sided 5\%,
Bonferroni over six rules), and the difference is also positive in two
alternative variants. The plan records that \ITM{10} was chosen as the headline
strike after the strike results had been seen.

\begin{table}[htbp]
\centering
\caption{Early exits under the calibrated model.}
\label{tab:exits}
\begin{tabular}{lrrrrrrr}
\toprule
Rule & sold early & mean & $t$ & profitable & 5th pct & worst & vs hold ($t$) \\
\midrule
H (hold) & 0\% & $+166.08$ & 4.25 & 54\% & $-1{,}186$ & $-3{,}601$ & --- \\
S3 & 51\% & $+86.76$ & 2.65 & 38\% & $-695$ & $-1{,}192$ & $-79.32$ ($-3.35$) \\
S5 & 35\% & $+111.02$ & 3.14 & 46\% & $-823$ & $-1{,}481$ & $-55.06$ ($-2.80$) \\
S8 & 20\% & $+136.54$ & 3.45 & 50\% & $-1{,}037$ & $-1{,}481$ & $-29.54$ ($-2.15$) \\
T5 & 44\% & $+83.79$ & 3.23 & 63\% & $-1{,}120$ & $-3{,}601$ & $-82.29$ ($-2.78$) \\
T10 & 19\% & $+101.36$ & 3.39 & 56\% & $-1{,}186$ & $-3{,}601$ & $-64.72$ ($-2.89$) \\
B5\_10 & 51\% & $+62.52$ & 2.34 & 48\% & $-810$ & $-1{,}481$ & $-103.56$ ($-3.79$) \\
\bottomrule
\end{tabular}
\tabnotes{\ITM{10}, variant (a), 644 modelled trades, level-normalised dollars per
contract, \$0.09 fee. ``Sold early'' is the share of trades the rule closes
before settlement. The last column is the mean paired difference from holding,
with its Newey--West $t$.}
\end{table}

\begin{figure}[htbp]
\centering
\caption{Early exits against holding to the close.}
\label{fig:exits}
\begin{tikzpicture}
\begin{axis}[paper, xbar, width=0.88\linewidth, height=4.9cm, bar width=8pt,
  xmin=0, xmax=195, ymajorgrids=false, xmajorgrids,
  ytick=data, yticklabels from table={figures/exits.csv}{label},
  y tick label style={font=\fontsize{10}{12}\selectfont}, enlarge y limits=0.09,
  xlabel={Mean P\&L per trade, \$ per contract (model, variant a, \ITM{10})},
  nodes near coords, nodes near coords align={horizontal}, point meta=x,
  every node near coord/.append style={font=\scriptsize, black!75,
    /pgf/number format/fixed, /pgf/number format/precision=0, /pgf/number format/print sign}]
\addplot[fill=series1, draw=none] table[x=per_trade, y=idx]{figures/exits.csv};
\end{axis}
\end{tikzpicture}
\tabnotes{644 modelled trades (2021--2026), \dollar0.09 fee. Every stop and
take-profit rule earned less on average than holding; stops bought a smaller
worst case (Table~\ref{tab:exits}).}
\end{figure}

Table~\ref{tab:exits} and Figure~\ref{fig:exits} show that no rule beat holding.
Every rule's mean difference from holding was negative, at every strike reported and
in every variant. Stops cut the worst loss by more than half, but they cost money
on average and reduced the share of profitable trades. The likely explanation is that the signal is a small drift over nine minutes,
so many adverse moves reverse before the close, and every early sale pays the
spread again. Take-profits raised the share of winning trades
and lowered the mean, because they cut the winners short and do nothing for the
losers. Sections~\ref{sec:fixedexit} and~\ref{sec:rtexit} revisit early exits
at real prices.

% =====================================================================
\section{Real option prices}
\label{sec:realprice}

The model studies leave two questions open, whether the edge is still there at
the prices actually quoted, and whether the current rule is the best nearby
setting. A pre-registered study addressed both with real OPRA quotes, tuning on
2013--2020 and confirming once on 2021--2026.

\subsection{Design}

At the entry minute, the spot $S$ is implied by put--call parity: it is the
median of $K + C_{\text{mid}} - P_{\text{mid}}$ over the three strikes where the
call and put mids are closest. Strikes then follow Section~\ref{sec:signal}. The
trade buys at the ask at the open of the entry minute and settles on the
official close, so its P\&L per contract is
\begin{equation}
\text{P\&L} = 100\,\bigl(\text{intrinsic value at the official close} - \text{ask}\bigr) - \text{fee}.
\end{equation}
A trade with no implied spot or no valid quote is dropped from that setting and
counted.

The search grid combines the candle minute (any one-minute candle from 15:40 to
15:50, entering at the next minute's open), the body threshold (0.40 to 0.90 in
steps of 0.05), the direction (with or against the candle) and the strike
(\ITM{20} to \OTM{20}), for 2{,}178 settings in all. A setting's score on
2013--2020 is its mean P\&L per day averaged with its neighbours at adjacent
thresholds, which favours a plateau over a sharp peak. A setting is eligible with
at least 150 trades, at least 50 in each half of the period, and a positive mean
in both halves. The best eligible setting would replace the current rule only if
its luck diagnostic gave $p \le 0.05$ and its per-day lead over the current rule
had $t \ge 2.0$. The chosen rule is confirmed on 2021--2026 if its mean P\&L per
day is positive with $t \ge 2.0$. The plan states that the 2021--2026 index
outcomes had already been used to test the signal, so the confirmation is fresh
for option prices but not for the index moves.

\subsection{Tuning on 2013--2020}

\begin{table}[htbp]
\centering
\caption{Tuning at real prices, 2013--2020: the eligible settings.}
\label{tab:rpdisc}
\setlength{\tabcolsep}{4pt}
\begin{tabular}{lrrrrrc}
\toprule
Cell (candle, body, direction, strike) & score & mean & trades & 2013--16 & 2017--20 & eligible \\
\midrule
15:42, $\ge 0.60$, with, \ATM & $+25.61$ & $+26.24$ & 250 & $+0.29$ & $+34.29$ & yes \\
15:50, $\ge 0.40$, with, \ATM & $+4.62$ & $+5.34$ & 302 & $+0.02$ & $+6.77$ & yes \\
\midrule
\emph{Current: 15:50, $\ge 0.60$, with, \ITM{10}} & $-12.65$ & $-7.70$ & 216 & $-21.57$ & $-3.82$ & no \\
\bottomrule
\end{tabular}
\tabnotes{2013-04-01 to 2020-12-31, real asks, \$0.09 fee, dollars per contract per
day. Only these two of the 2{,}178 settings were eligible. Luck diagnostic: 262
shifts, 142 of them with no eligible setting (counted as the worst); 95th
percentile of the shifted best score $+\$21.92$; $p = 0.034$. Best minus current:
$+\$40.17$ a day, Newey--West $t = 1.56$ over 593 days.}
\end{table}

Tuning on 2013--2020 had little to choose from, because only two of the 2{,}178
settings were eligible (Table~\ref{tab:rpdisc}). The better one, the 15:42 candle with a 60\% body at the money, passed the luck
diagnostic, but its lead over the current rule was not significant, so the
current rule went forward unchanged. The current rule itself lost money at the
ask in this period. Many of its signal days could not be priced at all, mostly
because too few strikes had valid quotes to imply a spot. In the confirmation
period only one signal day could not be priced.

\subsection{Confirmation on 2021--2026}

\begin{table}[htbp]
\centering
\caption{Confirmation: the current rule at real asks.}
\label{tab:rpconf}
\setlength{\tabcolsep}{4pt}
\begin{tabular}{lrrrrrrrr}
\toprule
Period & days & \$/day & $t$ & trades & \$/trade & won & median & worst \\
\midrule
\textbf{2021--2026} & \textbf{1{,}321} & $\mathbf{+29.87}$ & \textbf{2.40} & \textbf{585} & $\mathbf{+67.46}$ & \textbf{54\%} & $\mathbf{+44.91}$ & $\mathbf{-2{,}210}$ \\
2021--2022 & 376 & $+48.21$ & 2.01 & 157 & $+115.46$ & 57\% & $+65.91$ & $-1{,}361$ \\
2023--2026 & 945 & $+22.58$ & 1.56 & 428 & $+49.85$ & 53\% & $+35.91$ & $-2{,}210$ \\
\bottomrule
\end{tabular}
\tabnotes{The current rule (15:50 candle, body at least 60\%, with the candle,
\ITM{10}, held to the official close) at real asks, \$0.09 fee, dollars per
contract. $t$ is Newey--West on the per-day series. Subperiods are descriptive.
Mean per trade by year: 2021 $+\$175$ (67 trades), 2022 $+\$71$ (90), 2023
$+\$4$ (113), 2024 $+\$20$ (106), 2025 $+\$65$ (117), 2026 $+\$121$ (92).}
\end{table}

Table~\ref{tab:rpconf} shows that the current rule was confirmed. It earned $+\$29.87$ a
day with $t = 2.40$, against a bar of 2.0. The mean trade made $+\$67.46$, slightly
more than half of all trades were profitable, and the worst trade lost
\$2{,}210 on one contract. Figure~\ref{fig:rpyears} shows the full record. This confirmation has three weaknesses. It confirms the current rule at real prices but
does not show that it is the best setting, since no other setting was evaluated
on this period (Section~\ref{sec:refine} returns to that question). I had already seen the index outcomes. And the edge is weaker since 2023, when
it is not significant on its own, so much of the result comes from 2021--2022.

\begin{figure}[htbp]
\centering
\caption{The current rule at real option prices.}
\label{fig:rpyears}
\begin{tikzpicture}
\begin{axis}[paper, width=0.96\linewidth, height=5.6cm, xmin=2014, xmax=2027,
  xtick={2014,...,2026}, xticklabel style={/pgf/number format/1000 sep=},
  ymin=-20000, ymax=100000, ytick={-20000,0,...,100000},
  ylabel={Cumulative P\&L (\$)}, zeroline,
  title={(a) Running total, one contract per signal day}, title style={font=\fontsize{10}{12}\selectfont},
  legend style={at={(0.02,0.80)}, anchor=north west}]
\addplot[series1, line width=1.2pt] table[x=t, y=ask_cum]{figures/cumulative.csv};
\addlegendentry{Bought at the ask (as tested)}
\addplot[series2, line width=1pt, dash pattern=on 4pt off 2pt] table[x=t, y=mid_cum]{figures/cumulative.csv};
\addlegendentry{Bought at the mid (an upper bound; see Section~\ref{sec:fills})}
\draw[black!50, densely dotted, thick] (axis cs:2021,-20000) -- (axis cs:2021,100000);
\node[font=\scriptsize, black!70, anchor=north east] at (axis cs:2020.9,98000) {2013--2020: tuning};
\node[font=\scriptsize, black!70, anchor=north west] at (axis cs:2021.1,98000) {2021--2026: one confirmation test};
\end{axis}
\end{tikzpicture}

\vspace{2pt}
\begin{minipage}[t]{0.49\linewidth}\centering
\begin{tikzpicture}
\begin{axis}[paper, ybar, width=0.9\linewidth, height=5.0cm, bar width=6.5pt,
  xmin=2013.4, xmax=2026.6, xtick={2014,2016,...,2026}, minor xtick={2015,2017,...,2025},
  xticklabel style={/pgf/number format/1000 sep=},
  ymin=-8000, ymax=14000, ylabel={P\&L per year (\$)}, zeroline,
  title={(b) Total per calendar year, at the ask}, title style={font=\fontsize{10}{12}\selectfont}]
\addplot[fill=series1, draw=none] table[x=year, y=ask_total]{figures/by_year.csv};
\end{axis}
\end{tikzpicture}
\end{minipage}\hfill
\begin{minipage}[t]{0.49\linewidth}\centering
\begin{tikzpicture}
\begin{axis}[paper, width=0.93\linewidth, height=5.0cm,
  xmin=2013.4, xmax=2026.6, xtick={2014,2016,...,2026}, minor xtick={2015,2017,...,2025},
  xticklabel style={/pgf/number format/1000 sep=},
  ymin=0, ymax=4.5, ytick={0,1,...,4}, ylabel={Median full spread (\$ a share)},
  title={(c) Bid-ask spread of the contract bought}, title style={font=\fontsize{10}{12}\selectfont}]
\addplot[series1, line width=1.2pt, mark=*, mark size=1.6pt] table[x=year, y=median_spread]{figures/by_year.csv};
\end{axis}
\end{tikzpicture}
\end{minipage}
\tabnotes{The current rule (15:50 candle, body at least 60\%, with the candle,
\ITM{10}, held to the official close) at real Databento OPRA prices,
\dollar0.09 fee, dollars per contract at each year's prices. (a) The 801 priced
signal days, 2014--2026. At the ask the rule lost \dollar4{,}997 over the 216
trades of 2013--2020 and made \dollar39{,}463 over the 585 of 2021--2026. The
dashed line buys at the mid and is only an upper bound: a limit order at the mid
made about what paying the ask made, because it missed the best days
(Section~\ref{sec:fills}). (b) 2023 and 2024 were close to zero. (c) Median ask
minus bid at entry. The narrowing spread is why the same signal lost money at the
ask before 2021 and made money after it. Before 2016 SPXW had a same-day expiry
only on some days, so the early years have few trades. Sources: Databento OPRA;
Cboe; Dukascopy.}
\end{figure}

\subsection{The role of the bid-ask spread}
\label{sec:spread}

\begin{table}[htbp]
\centering
\caption{The current rule bought at the ask and at the mid.}
\label{tab:spread}
\begin{tabular}{llrrrrrc}
\toprule
Period & bought at & trades & \$/day & $t$ & \$/trade & won & median full spread \\
\midrule
\multirow{2}{*}{2013--2020} & ask & 216 & $-7.70$ & $-0.47$ & $-23.14$ & 42\% & \multirow{2}{*}{\dollar2.90} \\
 & mid & 216 & $+48.19$ & 2.87 & $+144.78$ & 56\% & \\
\midrule
\multirow{2}{*}{2021--2026} & ask & 585 & $+29.87$ & 2.40 & $+67.46$ & 54\% & \multirow{2}{*}{\dollar0.60} \\
 & mid & 585 & $+46.90$ & 3.71 & $+105.90$ & 56\% & \\
\bottomrule
\end{tabular}
\tabnotes{\ITM{10}, \$0.09 fee, dollars per contract. Rows at the mid are
descriptive and decide nothing. The spread is ask minus bid per share at entry,
the median over the priced trades. With a \$0.74 fee the ask rows become
$-\$23.79$ and $+\$66.81$ a trade ($t = 2.37$ in 2021--2026).}
\end{table}

Table~\ref{tab:spread} and Figure~\ref{fig:rpyears}(c) show how much the result depends on the spread. At the mid, the same trades made money in both periods. At the ask they
lost money in 2013--2020, when the median quoted spread was about five times its
2021--2026 level, and made money afterwards. Even since 2021, paying the ask instead of the mid costs about \$39 a trade,
more than a third of the edge at the mid. How the order is filled therefore
matters a great deal, and Section~\ref{sec:execution} looks at it.

\subsection{Strikes at real prices}

\begin{table}[htbp]
\centering
\caption{The current rule's signal at every strike, real asks, 2021--2026.}
\label{tab:rpstrikes}
\begin{tabular}{lrrrrrrr}
\toprule
Strike & trades & \$/trade & won & median & worst & \$/day & $t$ \\
\midrule
\ITM{20} & 586 & $+77.15$ & 54\% & $+51.91$ & $-2{,}670$ & $+34.20$ & 2.61 \\
\ITM{15} & 584 & $+71.00$ & 55\% & $+52.41$ & $-2{,}330$ & $+31.41$ & 2.43 \\
\textbf{\ITM{10}} & \textbf{585} & $\mathbf{+67.46}$ & \textbf{54\%} & $\mathbf{+44.91}$ & $\mathbf{-2{,}210}$ & $\mathbf{+29.87}$ & \textbf{2.40} \\
\ITM{5} & 585 & $+53.90$ & 50\% & $+2.91$ & $-1{,}900$ & $+23.87$ & 2.07 \\
\ATM & 583 & $+39.50$ & 39\% & $-110.09$ & $-1{,}640$ & $+17.46$ & 1.87 \\
\OTM{5} & 585 & $+12.18$ & 19\% & $-65.09$ & $-1{,}390$ & $+5.40$ & 0.78 \\
\OTM{10} & 581 & $+1.04$ & 11\% & $-25.09$ & $-1{,}180$ & $+0.46$ & 0.10 \\
\OTM{15} & 566 & $-6.96$ & 5\% & $-15.09$ & $-1{,}000$ & $-3.02$ & $-1.03$ \\
\OTM{20} & 529 & $-10.02$ & 3\% & $-10.09$ & $-840$ & $-4.19$ & $-2.34$ \\
\bottomrule
\end{tabular}
\tabnotes{\$0.09 fee, dollars per contract. Descriptive, not a second test:
choosing a strike from this table would be choosing on the confirmation period.}
\end{table}

\begin{figure}[htbp]
\centering
\caption{The current rule's signal at every strike, real asks, 2021--2026.}
\label{fig:rpstrikes}
\begin{minipage}[t]{0.49\linewidth}\centering
\begin{tikzpicture}
\begin{axis}[paper, ybar, width=0.95\linewidth, height=4.8cm, bar width=9pt,
  xtick=data, xticklabels from table={figures/strikes.csv}{label},
  xticklabel style={rotate=45, anchor=north east, font=\scriptsize}, enlarge x limits=0.07,
  ymin=-20, ymax=90, ytick={-20,0,...,80}, ylabel={Mean P\&L per trade (\$)}, zeroline,
  title={(a) Mean per trade}, title style={font=\fontsize{10}{12}\selectfont}]
\addplot[fill=series1, draw=none] table[x=idx, y=per_trade]{figures/strikes.csv};
\end{axis}
\end{tikzpicture}
\end{minipage}\hfill
\begin{minipage}[t]{0.49\linewidth}\centering
\begin{tikzpicture}
\begin{axis}[paper, ybar, width=0.95\linewidth, height=4.8cm, bar width=9pt,
  xtick=data, xticklabels from table={figures/strikes.csv}{label},
  xticklabel style={rotate=45, anchor=north east, font=\scriptsize}, enlarge x limits=0.07,
  ymin=0, ymax=60, ylabel={Trades that made money (\%)},
  title={(b) Share of trades that made money}, title style={font=\fontsize{10}{12}\selectfont}]
\addplot[fill=series1, draw=none] table[x=idx, y=won_pct]{figures/strikes.csv};
\end{axis}
\end{tikzpicture}
\end{minipage}
\tabnotes{\dollar0.09 fee, dollars per contract. In-the-money strikes earned the
most and won most often; out of the money the option rarely paid. Descriptive,
not a second test (Table~\ref{tab:rpstrikes}). Sources: Databento OPRA; Cboe;
Dukascopy.}
\end{figure}

Table~\ref{tab:rpstrikes} and Figure~\ref{fig:rpstrikes} show that at real prices
the strikes rank as they did in the model. Strikes 10 to 20 points in the money earned the most and won most often, and \ITM{10} needs about
half the premium of \ITM{20}. At the money the mean was positive but most trades
lost. Ten or more points out of the money, the option rarely paid.

\subsection{How accurate was the model?}
\label{sec:modelcheck}

\begin{table}[htbp]
\centering
\caption{The pricing model against real prices on the same trades.}
\label{tab:modelcheck}
\setlength{\tabcolsep}{4pt}
\begin{tabular}{lrrrrcrrr}
\toprule
& & \multicolumn{3}{c}{price per share} & ask / model & \multicolumn{3}{c}{P\&L per contract} \\
\cmidrule(lr){3-5}\cmidrule(lr){7-9}
Strike & $n$ & model & real mid & real ask & mean / median & model & real & real $-$ model ($t$) \\
\midrule
\ITM{20} & 586 & 20.09 & 20.07 & 20.58 & 1.02 / 1.02 & $+94.34$ & $+77.16$ & $-17.18$ ($-2.01$) \\
\ITM{15} & 584 & 15.37 & 15.22 & 15.68 & 1.02 / 1.01 & $+73.45$ & $+71.82$ & $-1.63$ ($-0.18$) \\
\textbf{\ITM{10}} & \textbf{585} & \textbf{10.29} & \textbf{10.51} & \textbf{10.89} & \textbf{1.06 / 1.05} & $\mathbf{+100.70}$ & $\mathbf{+67.77}$ & $\mathbf{-32.93}$ (\textbf{$-$3.87}) \\
\ITM{5} & 585 & 5.85 & 6.28 & 6.51 & 1.12 / 1.11 & $+96.90$ & $+51.57$ & $-45.33$ ($-5.63$) \\
\ATM & 583 & 2.40 & 3.08 & 3.18 & 1.50 / 1.33 & $+92.67$ & $+35.15$ & $-57.52$ ($-7.37$) \\
\OTM{5} & 585 & 0.93 & 1.37 & 1.44 & --- / 1.87 & $+45.59$ & $+9.32$ & $-36.27$ ($-5.52$) \\
\OTM{10} & 581 & 0.69 & 0.66 & 0.71 & --- / 2.65 & $-9.04$ & $+2.31$ & $+11.35$ (1.27) \\
\OTM{15} & 564 & 0.34 & 0.37 & 0.41 & --- / 15.31 & $-13.42$ & $-7.07$ & $+6.35$ (0.84) \\
\color{gray}\OTM{20} & \color{gray}524 & \color{gray}4.77 & \color{gray}0.23 & \color{gray}0.27 & \color{gray}--- / 10.82 & \color{gray}$-470.67$ & \color{gray}$-10.83$ & \color{gray}$+459.84$ (1.79) \\
\bottomrule
\end{tabular}
\tabnotes{The strike ladder's model (variant (a)) against real prices on the 644
strong-body trades of 2021--2026, at the index levels of the time (not rescaled),
on the days with a 15:51 quote for that contract. The model settles on the CFD's
15:59 close with a \$0.10 half-spread; the real trade buys at the ask and settles
on the official close; both pay \$0.09. Out of the money the mean ratio is not
meaningful, because model prices can be close to zero, so only the median is
shown.}
\end{table}

Table~\ref{tab:modelcheck} sets the model against the real trade on the same
days, contracts and index levels. On these terms the model implied
$+\$100.70$ a trade at \ITM{10} and the real trade made $+\$67.77$, a
significant shortfall. (The strike ladder's larger figures are in
level-normalised dollars.) In the money the model priced the option well, with
the real mid within about 2\% of the model price, but the real ask sat further
above the mid than the model's flat half-spread allowed. At the money the model
was too cheap, which is why its at-the-money results looked so much better than
the real ones. Far out of the money its skew factors made it unusable. The settlement price was not the cause, since the official close differed from the
CFD's 15:59 close by 1.33 points on average, with almost no bias.

% =====================================================================
\section{Execution}
\label{sec:execution}

The real-price results assume a fill at the quoted ask at 15:51:00. Here I ask whether that is realistic, whether a cheaper order would have done
better, and what a late entry or an early sale costs. Apart from the real-time
exit study of Section~\ref{sec:rtexit}, which was pre-registered, all of it is
descriptive. The limit-order and early-sale analyses were not pre-registered,
and 2021--2026 had already confirmed the rule. In brief, since 2021 the
contract traded near the ask on almost every signal day, neither limit orders
at the mid nor sales before the close improved on paying the ask and holding to
settlement, and later entries earned less.

\subsection{Did the contract trade near the ask?}
\label{sec:fillcheck}

\begin{table}[htbp]
\centering
\caption{Trading in the entry minute.}
\label{tab:fillcheck}
\begin{tabular}{lrr}
\toprule
Current rule, \ITM{10}, entry minute 15:51 & 2013--2020 & 2021--2026 \\
\midrule
Priced entries (all with a trade-bar file) & 216 & 585 \\
Traded in the entry minute & 122 (56\%) & 580 (99\%) \\
\quad median contracts traded in the minute & 10 & 75 \\
\quad lowest trade at or below the 15:51 ask & 95\% & 95\% \\
\quad lowest trade at or below the 15:51 mid & 52\% & 83\% \\
\quad minute's last trade minus the mid, mean & $+\dollar0.58$ & $+\dollar0.16$ \\
\quad minute's last trade minus the mid, median & $+\dollar0.57$ & $+\dollar0.25$ \\
\bottomrule
\end{tabular}
\tabnotes{From OPRA one-minute trade bars. Prices are per share (\$100 a contract
per \$1). The indented rows are shares or values over the entries that traded in
the minute.}
\end{table}

The ask is only a quoted price, and I do not observe what an order would
actually have paid. OPRA trade bars show
whether the contract traded in the entry minute and at what prices, and
Table~\ref{tab:fillcheck} summarises them. Since 2021 it traded in the entry minute on almost
every signal day, against just over half the days before, and on nearly all of
those days the minute's lowest trade was at or below the 15:51 ask. The minute's
last trade sat above the mid but inside the half-spread, and often above the
ask, as one would expect when the price keeps moving in the signal's direction.
A trade bar does not show the size available, the queue, or what a particular
order would have paid.

\subsection{Limit orders and late entry}
\label{sec:fills}

Since the trade made far more at the mid than at the ask
(Table~\ref{tab:spread}), the obvious thing to try is a limit order. I
simulated six order styles, fixed before any
result was computed: paying the 15:51 ask (P0); a limit at the 15:51 mid for one
minute, else no trade (P1); the same, but counting a trade at the limit as a fill
(P1t); P1 followed by paying the 15:52 ask if unfilled (P2); a mid limit left
working for three minutes (P3); and a limit halfway between the mid and the ask,
then the 15:52 ask (P4). A limit counts as filled if a trade printed strictly
below it during a minute it was working, or if the quote at the end of that
minute offered the option at or below it. These are proxies and not estimates of execution. Trade prints cannot show whether a resting order would have executed,
the order is assumed to be live with no delay, and one contract is assumed.

\begin{table}[htbp]
\centering
\caption{Order styles for the current rule, 2021--2026.}
\label{tab:fills}
\setlength{\tabcolsep}{4pt}
\begin{tabular}{lrrrrrr}
\toprule
& \multicolumn{4}{c}{2021--2026, 585 priced signal days} & \multicolumn{2}{c}{minus P0, \$/day} \\
\cmidrule(lr){2-5}\cmidrule(lr){6-7}
Order style & filled & \$/day & $t$ & minus P0 ($t$) & 2021--22 & 2023--26 \\
\midrule
P0 pay the 15:51 ask & 585 (100\%) & $+67.46$ & 2.29 & --- & --- & --- \\
P1 mid, 1 minute, else skip & 489 (84\%) & $+64.57$ & 2.37 & $-2.89$ ($-0.24$) & $-22.31$ & $+4.23$ \\
P1t as P1, touch counts & 501 (86\%) & $+61.04$ & 2.23 & $-6.42$ ($-0.56$) & $-22.31$ & $-0.59$ \\
P2 mid, 1 minute, then 15:52 ask & 585 (100\%) & $+69.63$ & 2.38 & $+2.17$ (0.55) & $+13.06$ & $-1.82$ \\
P3 mid, 3 minutes, else skip & 534 (91\%) & $+64.89$ & 2.31 & $-2.57$ ($-0.26$) & $+6.44$ & $-5.87$ \\
P4 halfway, 1 minute, then 15:52 ask & 585 (100\%) & $+62.23$ & 2.13 & $-5.23$ ($-1.63$) & $-2.36$ & $-6.29$ \\
\bottomrule
\end{tabular}
\tabnotes{\ITM{10}, \$0.09 fee, dollars per contract per priced signal day; an
unfilled order counts as \$0. The subperiods have 157 (2021--2022) and 428
(2023--2026) signal days. The $t$-statistics here treat signal days as
independent, so P0's $t$ is 2.29 rather than the per-day Newey--West 2.40 of
Table~\ref{tab:rpconf}.}
\end{table}

\begin{table}[htbp]
\centering
\caption{Why the mid limit does not help, and what a late entry costs.}
\label{tab:late}
\begin{tabular}{llrrr}
\multicolumn{5}{l}{\emph{(a) The same entries at the ask, split by whether P1's mid limit filled}} \\
\toprule
Period & P1 & days & \$/trade at the ask ($t$) & won \\
\midrule
2021--2026 & filled & 489 & $+42.52$ (1.31) & 52\% \\
 & \textbf{not filled} & \textbf{96} & $\mathbf{+194.49}$ \textbf{(2.80)} & \textbf{61\%} \\
\addlinespace
2013--2020 & filled & 69 & $+44.11$ (0.37) & 46\% \\
 & not filled & 147 & $-54.70$ ($-1.42$) & 40\% \\
\bottomrule
\end{tabular}

\bigskip
\begin{tabular}{lrrrr}
\multicolumn{5}{l}{\emph{(b) The same contract bought at a later minute's ask, \$/trade}} \\
\toprule
Buy at & 2013--2020 & 2021--2026 ($t$) & 2021--2022 & 2023--2026 \\
\midrule
15:51 & $-23.14$ & $+67.46$ (2.29) & $+115.46$ & $+49.85$ \\
15:52 & $-68.04$ & $+55.75$ (2.02) & $+111.07$ & $+35.46$ \\
15:53 & $-58.09$ & $+57.89$ (2.16) & $+128.30$ & $+32.06$ \\
15:55 & $-128.97$ & $+23.01$ (1.11) & $+51.10$ & $+12.71$ \\
\bottomrule
\end{tabular}
\tabnotes{Current rule, \ITM{10}, \$0.09 fee, dollars per contract. Panel (b) uses
the days with a quote at every one of these minutes: 216 (2013--2020), 585
(2021--2026), 157 (2021--2022) and 428 (2023--2026), which is every priced day.}
\end{table}

No order style beat paying the ask (Table~\ref{tab:fills}). On 2021--2026 every
style earned roughly what paying the ask earned, and every difference was
within noise. Panel (a) of Table~\ref{tab:late} suggests that the cause is
adverse selection. A limit at the mid fills when someone sells below the mid, which
happens most when the option is getting cheaper, that is, when the trade is
going wrong. When the trade works at once, the price moves away and the limit is
left unfilled. The days on which the mid limit filled earned well below the
average at the ask, and the days it missed earned far more. The limit saves money
on the days it fills and misses the best days, and the two effects roughly
cancel. Before 2021, with spreads about five times wider, the mid limit did help,
but on a small sample.

Entering late also costs money. Bought at a later minute's ask, the same
contract earned less the later the entry (Table~\ref{tab:late}, panel b), about
\$12 a trade less at 15:52 and about two-thirds less at 15:55. The rule's
results assume entry at 15:51:00 exactly.

\subsection{Selling before the close at real bids}
\label{sec:fixedexit}

As a check on the model's exit results, each of the 585 confirmation trades,
bought at the 15:51 ask, was also sold at the bid at the start of each later
minute, paying the fee twice. Every fixed exit minute from 15:52 to 15:59 earned
less than holding to settlement, by \$79 to \$136 a trade, and no year favoured
an early exit. Selling at the best minute in hindsight would have earned far
more, which shows only how much the price moves after entry.

\subsection{Real-time exit rules at real bids}
\label{sec:rtexit}

A fixed exit minute ignores what the price does after entry. A pre-registered
study tested rules that watch the option's own bid and sell when it moves. Each of
the rule's 270 entries of 2021--2023, bought at the 15:51 ask at \ITM{10}, was
followed through one-minute quote snapshots to 15:59. A rule decides at each
minute from 15:52 to 15:58 and sells at the next minute's bid, paying a second
fee; a trade it never sells settles. Twenty-six rules were fixed in advance:
trailing stops that sell when the bid falls 10, 20, 30, 40 or 50\% below its
highest point since entry (5 rules); profit locks that, once the bid has risen
10, 20 or 30\% above the price paid, sell if it gives back a third, two-thirds or
all of that best gain (9); and take-profits that sell once the bid is 10, 20 or
30\% above the price paid, from 15:52, 15:54, 15:56 or 15:58 (12). A rule was
eligible if it sold early on at least 25 trades, at least 10 in each half of the
tuning period (January 2021 to September 2022 and October 2022 to December
2023), and beat holding on average in both halves. The eligible rule with the
best score, its mean improvement on holding averaged with those of its
neighbours in the grid, would then have faced a block sign-flip luck diagnostic
($p \le 0.05$) and a single test on 2024--2026, passing if its mean improvement
was positive with Newey--West $t \ge 2.0$.

The study ran under two plans. The first (locked 2026-11-10 22:17 UTC) also
included 18 rules that trail the index level implied by put--call parity. It
required complete quotes for at least 95\% of the entries and stopped at that
gate before any rule was computed: only 197 of the 270 entries were complete,
because the implied index level, which needs three strikes with a bid on both
sides, was often missing in the last minutes (on 70 entries at 15:58 alone). A revised plan (locked 2026-11-11 00:01 UTC), written after that stop
and its counts were known, kept only the 26 rules that read the option's own bid
and required complete option quotes, which all 270 entries had. It was run once.

\begin{table}[htbp]
\centering
\caption{Real-time exit rules at real bids.}
\label{tab:rtexit}
\begin{tabular}{lrrlr}
\toprule
 & & \multicolumn{2}{c}{2021--2023 (tuning)} & 2024--2026 \\
\cmidrule(lr){3-4}
Family & Rules & Against holding & Least bad & Against holding \\
\midrule
Trailing stop & 5 & $-120.44$ to $-13.79$ & TS50 & $-30.84$ to $+3.97$ \\
Profit lock & 9 & $-37.86$ to $-8.94$ & PL30\_100 & $-26.96$ to $+8.62$ \\
Take-profit & 12 & $-106.88$ to $-49.12$ & TP1558\_30 & $-108.42$ to $-39.28$ \\
Sale at a fixed minute & 8 & $-168.26$ to $-111.74$ & 15:57 & $-120.04$ to $-51.12$ \\
\bottomrule
\end{tabular}
\tabnotes{\ITM{10}, bought at the 15:51 ask; dollars per contract. ``Against
holding'' is the range, across the family's rules, of the mean paired difference
from holding to settlement over all trades, counting zero where a rule held.
Holding earned $+\$68.78$ a trade on 2021--2023 (270 trades) and $+\$66.32$ on
2024--2026 (315). TS50 sells when the bid falls 50\% below its peak; PL30\_100,
once the bid has been 30\% above the price paid, sells if it falls back to that
price; TP1558\_30 sells from 15:58 once the bid is 30\% above the price paid. The
2024--2026 column was computed after the decision and decides nothing; the
fixed-minute row is descriptive in both periods.}
\end{table}

Table~\ref{tab:rtexit} reports the outcome. No rule was eligible, so neither the
luck diagnostic nor the test was run. Every one of the 26 rules earned less than
holding on 2021--2023, by \$9 to \$120 a trade. Take-profits did worst. Selling once the bid was 10 to 30\% above the price paid cost \$49 to \$107 a trade
against holding, with Newey--West $t$ between $-3.2$ and $-5.6$. They sell the trades that would have kept rising and do nothing for the losers,
and every early sale gets the bid while holding settles at the index value with
no spread. I looked at 2024--2026 after the decision and it agrees:
take-profits trailed
holding by \$39 to \$108 a trade, and the trailing stops and profit locks were
within \$31 of holding, none with $|t|$ above 1.6. Selling at the best minute in
hindsight would have added \$224 a trade on 2021--2023, which again measures how much the price moves and is not a gain any rule here
could capture.

% =====================================================================
\section{Re-tuning on recent data}
\label{sec:refine}

Tuning on 2013--2020 could not tell nearby thresholds or candle minutes apart,
because spreads were wide, many signal days could not be priced, and only two
settings were eligible. A further pre-registered study repeated the search on
2021--2023, the first three years with tight spreads and, from 2022, daily SPXW
expiries. It would test its pick once on 2024--2026. The pick failed the luck
diagnostic, so nothing was tested and the current rule was kept.

Pricing and the grid of 2{,}178 settings are as in
Section~\ref{sec:realprice}. A setting is eligible with at least 100 trades, at
least 40 in each half of 2021--2023, a positive mean in each half, and at least
95\% of its signal days priced. The pick goes to the test only if it lies
outside the current rule's family (the 15:50 candle, a 60\% body and the
candle's direction, at any strike), its luck diagnostic gives $p \le 0.05$, and
its paired difference from the current rule on 2021--2023 is positive. On
2024--2026 the pick would be adopted only if both its mean per day and its
paired lead over the current rule were positive with $t \ge 2.0$. Tuning and
testing were separate runs: the tuning step could not read any date after 2023,
and it recorded its decision, with a hash, before the test step could run.

\begin{table}[htbp]
\centering
\caption{The re-tuning pick and the current rule.}
\label{tab:refine}
\setlength{\tabcolsep}{4pt}
\begin{tabular}{llrrrrrrr}
\toprule
Rule & period & days & \$/day & $t$ & trades & \$/trade & won & worst \\
\midrule
\multirow{3}{*}{\shortstack[l]{Pick: 15:50, $\ge 0.80$,\\ with, \ITM{20}}}
 & \textbf{2021--2023} & \textbf{625} & $\mathbf{+39.70}$ & \textbf{2.94} & \textbf{133} & $\mathbf{+186.55}$ & \textbf{62\%} & $\mathbf{-1{,}687}$ \\
 & H1 & 315 & $+75.61$ & 3.60 & 70 & $+340.22$ & 66\% & $-1{,}687$ \\
 & H2 & 310 & $+3.21$ & 0.21 & 63 & $+15.80$ & 59\% & $-1{,}433$ \\
\multirow{3}{*}{\shortstack[l]{Current: 15:50, $\ge 0.60$,\\ with, \ITM{10}}}
 & 2021--2023 & 624 & $+29.76$ & 1.78 & 270 & $+68.78$ & 55\% & $-1{,}361$ \\
 & H1 & 314 & $+63.52$ & 2.31 & 133 & $+149.97$ & 56\% & $-1{,}361$ \\
 & H2 & 310 & $-4.43$ & $-0.25$ & 137 & $-10.03$ & 54\% & $-1{,}340$ \\
\midrule
\multicolumn{9}{l}{\emph{2024-01-01 to 2026-10-20, computed after the decision was written (descriptive)}} \\
Pick & 2024--2026 & 697 & $+18.88$ & 1.30 & 167 & $+78.80$ & 51\% & $-2{,}670$ \\
Current & 2024--2026 & 697 & $+29.97$ & 1.64 & 315 & $+66.32$ & 53\% & $-2{,}210$ \\
\bottomrule
\end{tabular}
\tabnotes{Real asks, official close, \$0.09 fee, dollars per contract. $t$ is
Newey--West on the per-day series. H1 is 2021-01-01 to 2022-09-30 and H2 is
2022-10-01 to 2023-12-31. The pick failed the luck diagnostic, so the 2024--2026
rows decide nothing; the current rule's results there were already known.}
\end{table}

Table~\ref{tab:refine} reports the outcome. On 2021--2023, 119 settings were
eligible. The best was the 15:50 candle with a body of at least 80\%, traded with
the candle 20 points in the money. It failed the luck diagnostic ($p = 0.21$):
about one shift in five found a best score at least as good with signals
unrelated to the days. No candidate advanced, the test on 2024--2026 was not run, and the current rule
was kept. The pick's tuning record came almost entirely
from the first half of the period.

Computed only after the decision was recorded, and deciding nothing, the pick
earned less than the current rule on 2024--2026, as one would expect of a winner chosen by luck. Figure~\ref{fig:threshold} shows why the threshold is hard to tune. On 2021--2023 a larger body earned more per trade
on far fewer trades, but the pattern did not persist; per day, the current 60\%
threshold earned about the same in both periods, and no threshold beat it in
both. Among the other candle minutes, the best eligible 15:48 setting was
steadier than the pick across the two halves, but it is one observation from a
large search and remains untested. Even with tight spreads, three years of data
cannot tell the nearby settings apart, and choosing the best past result would
most likely choose a lucky one.

\begin{figure}[htbp]
\centering
\caption{The body threshold.}
\label{fig:threshold}
\begin{minipage}[t]{0.49\linewidth}\centering
\begin{tikzpicture}
\begin{axis}[paper, width=\linewidth, height=5.0cm, xmin=0.38, xmax=0.92,
  xtick={0.4,0.5,...,0.9}, xticklabel style={/pgf/number format/fixed,
    /pgf/number format/precision=1, /pgf/number format/fixed zerofill},
  ymin=0, ymax=45, ylabel={Mean P\&L per day (\$)},
  xlabel={Body at least (share of the candle's range)},
  title={(a) Per day}, title style={font=\fontsize{10}{12}\selectfont},
  legend to name=thrlegend, legend columns=2,
  legend style={/tikz/every even column/.append style={column sep=1.2em}}]
\addplot[series1, line width=1.2pt, mark=*, mark size=1.6pt] table[x=theta, y=tune_per_day]{figures/threshold.csv};
\addlegendentry{2021--2023 (tuning)}
\addplot[series2, line width=1.1pt, dash pattern=on 4pt off 2pt, mark=square*, mark size=1.5pt,
  mark options={solid}] table[x=theta, y=test_per_day]{figures/threshold.csv};
\addlegendentry{2024--2026 (looked at only after the decision)}
\draw[black!50, densely dotted, thick] (axis cs:0.6,0) -- (axis cs:0.6,45);
\node[font=\scriptsize, black!70, anchor=south west] at (axis cs:0.605,1) {current};
\end{axis}
\end{tikzpicture}
\end{minipage}\hfill
\begin{minipage}[t]{0.49\linewidth}\centering
\begin{tikzpicture}
\begin{axis}[paper, width=\linewidth, height=5.0cm, xmin=0.38, xmax=0.92,
  xtick={0.4,0.5,...,0.9}, xticklabel style={/pgf/number format/fixed,
    /pgf/number format/precision=1, /pgf/number format/fixed zerofill},
  ymin=0, ymax=200, ylabel={Mean P\&L per trade (\$)},
  xlabel={Body at least (share of the candle's range)},
  title={(b) Per trade}, title style={font=\fontsize{10}{12}\selectfont}]
\addplot[series1, line width=1.2pt, mark=*, mark size=1.6pt] table[x=theta, y=tune_per_trade]{figures/threshold.csv};
\addplot[series2, line width=1.1pt, dash pattern=on 4pt off 2pt, mark=square*, mark size=1.5pt,
  mark options={solid}] table[x=theta, y=test_per_trade]{figures/threshold.csv};
\draw[black!50, densely dotted, thick] (axis cs:0.6,0) -- (axis cs:0.6,200);
\node[font=\scriptsize, black!70, anchor=south west] at (axis cs:0.605,4) {current};
\end{axis}
\end{tikzpicture}
\end{minipage}

\vspace{2pt}
\ref{thrlegend}
\tabnotes{15:50 candle, with the candle, \ITM{10}, real asks, \dollar0.09 fee,
dollars per contract. On 2021--2023 a bigger body earned more per trade, on far
fewer trades (407 at 0.40, 56 at 0.90). The pattern did not hold on 2024--2026.
Per day, 0.60 made about \dollar30 in both periods; the 2021--2023 favourite,
0.80, made \dollar13 a day afterwards (Table~\ref{tab:threshold}). Sources:
Databento OPRA; Cboe; Dukascopy.}
\end{figure}

% =====================================================================
\section{Intraday momentum and dealer gamma}
\label{sec:momentum}



\citet{ghlz2018} find that the S\&P~500's first half-hour return, measured from
the previous close, predicts its last half hour on SPY from 1993 to 2013.
\citet{bdlm2021} find that the return from the previous close to 15:30 does the
same across futures markets from 1974 to 2020, and attribute the effect to option
dealers hedging when they are net short gamma. The 15:50 candle sits inside that
last half hour, so momentum might improve the candle trade, or the candle might be momentum in
another form. A pre-registered study asked whether the momentum appears in my
data (Part A), whether it improves the candle trade (Part B), and whether
either is stronger when dealers are short gamma (Part C).

With PC the previous session's close and prices taken at bar
opens, the predictor is $r_{\text{rod}} = \ln(P_{15:30}/\mathrm{PC})$ and the
target is $r_{\text{last}} = \ln(X/P_{15:30})$, where $X$ is the close of the
15:59 bar. Part A regresses $r_{\text{last}}$ on $r_{\text{rod}}$ over 2013--2026
and requires a positive slope with $t \ge 2.0$. It is an adapted replication: it
uses log returns on a CFD rather than simple returns on SPY or futures, and only
2021--2026 lies outside both papers' samples. Part B tests four filters that keep
the candle trade only when it agrees with the day's earlier move. A filter must
improve the trade with $p \le 0.0125$ on 2013--2020, and also in both halves and
in a version that uses index returns instead of option prices. Part C builds a
net-gamma proxy from the open interest of every SPX and SPXW option, following
\citet{bdlm2021} and assuming that dealers are long every call and short every
put. It asks whether momentum (C1) or the candle trade (C2) is stronger when the
proxy is negative, each with $p \le 0.025$.

\begin{table}[htbp]
\centering
\caption{Part A: the momentum regression.}
\label{tab:momA}
\setlength{\tabcolsep}{4pt}
\begin{tabular}{llrrrr}
\toprule
Regression of $r_{\text{last}}$ on & sample & $n$ & slope $\times 100$ & $t$ & $R^2$ \\
\midrule
\textbf{$r_{\text{rod}}$ (primary)} & \textbf{2013--2026} & \textbf{3{,}238} & $\mathbf{+1.567}$ & \textbf{1.02} & \textbf{0.0030} \\
$r_{\text{rod}}$ & 2013--2020 & 1{,}802 & $+2.866$ & 1.16 & 0.0086 \\
$r_{\text{rod}}$ & 2021--2026 & 1{,}436 & $-0.173$ & $-0.18$ & 0.0000 \\
$r_{\text{rod}}$, simple returns & 2013--2026 & 3{,}238 & $+1.587$ & 1.04 & 0.0030 \\
$r_{\text{rod}}$, winsorised 1\%/99\% & 2013--2026 & 3{,}238 & $+1.189$ & 1.08 & 0.0013 \\
$r_{\text{first}}$ (first half hour) & 2013--2026 & 3{,}236 & $+1.517$ & 0.57 & 0.0014 \\
$r_{12}$ (15:00--15:30) & 2013--2026 & 3{,}204 & $+9.510$ & 1.30 & 0.0043 \\
\midrule
\emph{BDLM, E-mini futures} & \emph{to May 2020} & & \emph{about 6.2} & \emph{about 5.0} & \emph{about 0.034} \\
\bottomrule
\end{tabular}
\tabnotes{An adapted replication on the CFD, 2013-04-01 to 2026-10-20. Only the
first row is a test; the rest are descriptive. BDLM is \citet{bdlm2021}. By
single year, the $t$ of the $r_{\text{rod}}$ slope was 2.12, 1.75 and 1.84 in
2013--2015, and between $-1.85$ and $+0.96$ in every year since.}
\end{table}

The momentum did not replicate (Table~\ref{tab:momA}). The slope was positive but insignificant over 2013--2026,
mildly positive only in 2013--2015, and about zero on 2021--2026, the only period
outside both papers' samples. This does not show that the published results are
wrong, since samples, instruments and periods differ. It shows only that the
effect is not visible in these data since 2013. No momentum filter improved the
candle trade either (Table~\ref{tab:momB}, Appendix~\ref{app:tables}): trades that agreed with the day's move did better
per trade, but the trades a filter skipped were still profitable on average, so
skipping them lowered the total. Neither gamma test passed either
(Table~\ref{tab:partc}, Appendix~\ref{app:tables}). Under model prices the candle trade did better on
negative-gamma days, but the difference missed its bar ($p = 0.029$ against
0.025), came mostly from 2013--2020, and was absent at real asks on 2021--2026.
The proxy therefore gives no reason to skip or add trades. Its positioning is assumed and not observed, and open interest misses positions opened and
closed on the same day, which is most 0DTE trading since 2022.

% =====================================================================
\section{Smaller contracts: XSP and SPY}
\label{sec:mini}


The current rule's SPX contract cost about \$1{,}100 on 2021--2026, more than a
small account holds. Two underlyings have 0DTE options about a tenth the size:
XSP (Mini-SPX, cash-settled, one tenth of SPX) and SPY (the ETF, settled in
shares). A pre-registered study traded the same signal on both at real quotes.
On each, the spot is implied from that root's own 15:51 quotes, and the contract
is \$1 in the money in the candle's direction (about 10 SPX points), bought at
the 15:51 ask. XSP is held to the official SPX close divided by ten. SPY, which
settles in shares, is sold at the 15:59 bid. Each is confirmed only if its mean
P\&L per day is positive with $t \ge 2.0$ on 2021--2026. Because the SPX outcomes
on these days were already known, the study measures the cost of trading the
signal on smaller contracts, not whether the signal works.

\begin{table}[htbp]
\centering
\caption{XSP and SPY at real quotes, 2021--2026.}
\label{tab:mini}
\begin{tabular}{@{}lrrrrrrr@{}}
\toprule
Contract & Days & \$/day & $t$ & Trades & \$/trade & Won & Per \$1 \\
\midrule
XSP \ITM{1}, held & 1{,}025 & $-15.16$ & $-5.48$ & 313 & $-49.66$ & 35\% & $-0.285$ \\
SPY \ITM{1}, sold 15:59 bid & 1{,}262 & $+0.23$ & $+0.23$ & 560 & $+0.51$ & 45\% & $+0.004$ \\
\midrule
SPX \ITM{10}, held & 1{,}321 & $+29.87$ & $+2.40$ & 585 & $+67.46$ & 54\% & $+0.061$ \\
\bottomrule
\end{tabular}
\tabnotes{One contract per signal. ``Per \$1'' is total P\&L over total premium
paid. The SPX row rebuilds the confirmation of Table~\ref{tab:rpconf} from its
trade list.}
\end{table}

Neither was confirmed (Table~\ref{tab:mini}). XSP lost money reliably and SPY
broke even, and the reason in both cases is cost. The SPX trade's edge was about six cents per dollar
of premium, and both smaller contracts cost more than that to trade. XSP's
median quoted spread was 31\% of the option's price, against 5\% for SPX on the
same days. It has narrowed since 2021 and XSP's results improved with it, but
even in 2026 the trade only about broke even. XSP is also thin: on almost half
of the signal days too few strikes had two-sided quotes to imply a spot. SPY's
spread is small, but a round trip pays it twice and gives up the final minute.
For XSP, buying at the mid, other strikes, a higher fee and a cap on the premium
all left the conclusion unchanged. SPY traded at the mid on both sides would have
been profitable, but that is a price an order cannot count on
(Section~\ref{sec:fills}).

A vertical debit spread does not need a separate test. Bought
at the long leg's ask and sold at the short leg's bid, it earns exactly the long
leg's P\&L at the ask, less the short leg's P\&L at the ask, less the short
leg's quoted spread and a fee, and every term but the spread is already in
Table~\ref{tab:rpstrikes}. With average spreads, no vertical did better per
dollar of premium than its long leg alone.

% =====================================================================
\section{An independent re-validation}
\label{sec:audit}

After the studies were complete, the whole chain was rebuilt with new code that
shares nothing with the study scripts and reads the same raw files. No new data
were bought. The check is descriptive, because every day it uses had already been used. It
can verify the earlier work and measure how much the data can say, but it
cannot test a new setting. The re-validation reproduced the results, placed the
start of the effect near a 60\% body, and showed that the CFD candle from which
the signals were computed differs in shape from the SPX index candle.

The first check was the integrity of the data. The index bars contain no duplicate, out-of-order or
internally inconsistent bar, and every kept day has option quotes, or a record
that no same-day expiry existed, and an official close. The weekdays without an
official close are exactly the exchange holidays, the official closes equal the
broker's daily closes on every day that has both, and every quote snapshot is
stamped exactly on the minute, with no duplicate contract. At 15:51 the CFD price
lies within about half an index point, on average, of the spot implied by
put--call parity from the real quotes. The 15:50 bar is also the right minute:
its body moves almost one for one with the option-implied change over the same
minute, and not with the minutes before or after. One weakness was found. The
CFD feed was coarser in its early years, when between a fifth and a third of
15:50 bars had a body equal to the whole range, so a 60\% threshold was a looser
filter in the early discovery data than it is now.

Rebuilt from the raw quotes, the confirmation of
Section~\ref{sec:realprice} is identical trade for trade, and so are the
2013--2020 trades, the strike ladder, the yearly means, the late entries, the
fixed-minute exits and the threshold profile of the re-tuning study. Coverage
differs sharply between the two periods. Before 2021 the implied spot, which
needs three strikes quoted on both sides, was often unavailable on calm days, so
only a little over half of the signal days could be priced, against all but one
since 2021. The missing days do not hide a better result: priced with the CFD
price as the spot, as a sensitivity, they lost money at the ask as well
($-\$89$ a trade, against $+\$68$ at the mid). A five-day block bootstrap of
2021--2026 puts the current rule's mean between $+\$21$ and $+\$115$ a trade
(90\% interval), so the mean is positive but imprecisely measured.

\begin{table}[htbp]
\centering
\caption{The 15:50 candle by body-to-range bucket.}
\label{tab:auditbuckets}
\begin{tabular}{@{}lrrrrrrr@{}}
\toprule
 & \multicolumn{2}{c}{Index, 2013--2020} & \multicolumn{2}{c}{Index, 2021--2026} & \multicolumn{3}{c}{Options, 2021--2026} \\
\cmidrule(lr){2-3}\cmidrule(lr){4-5}\cmidrule(lr){6-8}
Body / range & bps & $t$ & bps & $t$ & Trades & Ask & Mid \\
\midrule
0.0--0.2 & $-1.46$ & $-0.76$ & $-0.29$ & $-0.29$ & 220 & $-53$ & $-13$ \\
0.2--0.4 & $+0.20$ & $+0.23$ & $-0.22$ & $-0.26$ & 223 & $-21$ & $+18$ \\
0.4--0.5 & $+1.38$ & $+1.34$ & $-0.16$ & $-0.15$ & 141 & $-15$ & $+21$ \\
0.5--0.6 & $-2.35$ & $-1.22$ & $-0.24$ & $-0.20$ & 149 & $-59$ & $-17$ \\
0.6--0.7 & $+3.97$ & $+2.25$ & $+1.71$ & $+1.49$ & 150 & $+6$ & $+46$ \\
0.7--0.8 & $+2.77$ & $+2.75$ & $+1.26$ & $+1.30$ & 135 & $+42$ & $+81$ \\
0.8--0.9 & $+0.73$ & $+0.70$ & $+3.50$ & $+2.40$ & 149 & $+97$ & $+138$ \\
0.9--1.0 & $+2.17$ & $+2.48$ & $+2.92$ & $+2.61$ & 151 & $+122$ & $+156$ \\
\midrule
Below 0.6 & $-0.54$ & $-0.71$ & $-0.23$ & $-0.46$ & 733 & $-37$ & $+2$ \\
At least 0.6 & $+2.42$ & $+3.61$ & $+2.36$ & $+4.25$ & 585 & $+67$ & $+106$ \\
\bottomrule
\end{tabular}
\tabnotes{Index move: CFD, 15:51 open to 15:59 close, in the candle's direction, in basis
points, with a Newey--West (lag 5) $t$-statistic on that bucket's days. Option
columns: \ITM{10} in the candle's direction, 2021--2026, dollars per trade at the
15:51 ask and at the mid. One trade lies exactly on the 0.80 boundary; classified
in exact rather than floating-point arithmetic, it moves the option means at the
ask of the 0.7--0.8 and 0.8--0.9 buckets to $+\$45$ and $+\$94$. Descriptive.}
\end{table}

To see where the effect begins, Table~\ref{tab:auditbuckets} splits the 15:50
candle by the size of its body. Below 60\% there is no effect in either
period, in the index or in the option trade. From 60\% up, the index moved in
the candle's direction by about the same amount in both periods. Within the
signal, larger bodies earned more per trade in 2021--2026, but no single option
bucket is significant on its own, and the ordering is not stable: the 60--70\% bucket,
about break-even at the ask in 2021--2026, was the strongest in 2013--2020.

\begin{table}[htbp]
\centering
\caption{The body threshold at real asks, 2021--2026.}
\label{tab:auditgrid}
\begin{tabular}{@{}lrrrrr@{}}
\toprule
Body at least & Trades & \$/trade & \$/day & $t$ & Best \\
\midrule
0.50 & 734 & $+42$ & $+23.3$ & 1.70 & 2\% \\
0.55 & 656 & $+47$ & $+23.2$ & 1.80 & 1\% \\
0.60 (current) & 585 & $+67$ & $+29.9$ & 2.40 & 32\% \\
0.65 & 508 & $+68$ & $+26.3$ & 2.27 & 2\% \\
0.70 & 435 & $+89$ & $+29.2$ & 2.65 & 22\% \\
0.75 & 380 & $+92$ & $+26.4$ & 2.56 & 7\% \\
0.80 & 300 & $+110$ & $+24.9$ & 2.58 & 8\% \\
0.85 & 219 & $+99$ & $+16.5$ & 1.94 & 1\% \\
0.90 & 151 & $+122$ & $+13.9$ & 2.09 & 1\% \\
\bottomrule
\end{tabular}
\tabnotes{15:50 candle, with the candle, \ITM{10} at the ask, 2021-01-01 to 2026-10-20.
``Per day'' means per session with a same-day expiry and prices (1{,}321 of
1{,}446 sessions). ``Best'' is the share of 5{,}000 five-day block-bootstrap
resamples in which that threshold had the highest mean per day (grid 0.00 to 0.95
in steps of 0.05); the middle 90\% of the resampled optimum ran from 0.25 to
0.80, and the 90\% interval for 0.70 minus 0.60 is $-\$10.7$ to $+\$8.8$ a day.
A walk-forward choice made $+\$16$ a day on 2022--2026 when its history starts in
2021 and $+\$20$ when it starts in 2013, against $+\$24$ for a fixed 0.60
($t = -0.98$ and $-0.96$ for the differences). 0.80 made $+\$38$ a day on
2021--2023 and $+\$13$ on 2024--2026. Descriptive.}
\end{table}

Above 60\%, the thresholds cannot be ranked. Raising the threshold trades
frequency for size, as Table~\ref{tab:auditgrid} shows. Per day, the difference
between 60\% and 70\% is statistically unresolved: no threshold differs
significantly from 60\% on the paired daily series, and the bootstrap interval
for the difference between 0.70 and 0.60 straddles zero. The bootstrap shares in
the table describe how unstable the in-sample optimum is; they are not probabilities that a given threshold is the best. A walk-forward choice, re-picking
each year the threshold that had been best on all earlier days, earned less than
a fixed 0.60 whichever history it started from, and with the longer history it
chose 0.60 itself in most years. The two halves of 2021--2026 also disagree above
0.75. None of this is a fresh test, since the same days chose the rule. Within
this sample the effect starts near 60\%, and the thresholds from 60\% to about
75\% cannot be ranked.

The body-to-range ratio works as a cut-off and does not measure the strength of
the signal. The bucket tables show where the effect sits, and correlations show
how much of each day it explains (Table~\ref{tab:auditcorr}). Across all days a higher body-to-range
ratio went with a larger index move in the candle's direction and a better
option result, significantly but faintly: the largest of these correlations
accounts for 0.7\% of the day-to-day variation. Among signal days alone the
relation is not distinguishable from zero in either period, and the thirds of
the signal days by ratio show no stable slope (Table~\ref{tab:auditthirds}): in
2013--2020 the lowest third had the largest index move, and in 2021--2026 it had
the smallest option result. Once a candle passes the threshold, a larger ratio
does not mean a larger gain, which fits the finding that the thresholds above
60\% cannot be ranked.

\begin{table}[htbp]
\centering
\caption{Correlations between the 15:50 CFD candle and what followed.}
\label{tab:auditcorr}
\begin{tabular}{@{}llllrrrrr@{}}
\toprule
Feature & Outcome & Days & Period & $n$ & $r$ & $p$ & $\rho$ & $p$ \\
\midrule
Body / range & Index move & All & 2013--2020 & 1{,}785 & $+0.06$ & 0.007 & $+0.04$ & 0.066 \\
Body / range & Index move & All & 2021--2026 & 1{,}445 & $+0.08$ & 0.002 & $+0.07$ & 0.011 \\
Body / range & Option result & All & 2021--2026 & 1{,}318 & $+0.08$ & 0.005 & $+0.07$ & 0.018 \\
\midrule
Body / range & Index move & Signal & 2013--2020 & 908 & $-0.05$ & 0.108 & $-0.04$ & 0.190 \\
Body / range & Index move & Signal & 2021--2026 & 644 & $+0.04$ & 0.287 & $+0.02$ & 0.605 \\
Body / range & Option result & Signal & 2021--2026 & 585 & $+0.06$ & 0.177 & $+0.05$ & 0.193 \\
\midrule
Body size & Index move & Signal & 2013--2020 & 908 & $+0.19$ & $<0.001$ & $+0.08$ & 0.023 \\
Body size & Index move & Signal & 2021--2026 & 644 & $+0.13$ & 0.001 & $+0.11$ & 0.006 \\
Body size & Right direction & Signal & 2013--2020 & 901 & $+0.02$ & 0.617 & $+0.03$ & 0.438 \\
Body size & Right direction & Signal & 2021--2026 & 644 & $+0.04$ & 0.255 & $+0.04$ & 0.367 \\
Body size & Option result & Signal & 2021--2026 & 585 & $+0.12$ & 0.003 & $+0.08$ & 0.055 \\
\bottomrule
\end{tabular}
\tabnotes{Body size is $|c-o|/o$. Index move: CFD, 15:51 open to 15:59 close, in
the candle's direction. Option result: \ITM{10} in the candle's direction at the
15:51 ask, held to settlement. ``Right direction'' is 1 if the index moved the
candle's way. Signal days have a body of at least 60\% of the range; days with a
flat 15:50 candle are excluded. Periods: 2013-04-01 to 2020-12-31 and 2021-01-01
to 2026-10-20. $r$ is Pearson's and $\rho$ Spearman's coefficient; $p$-values
assume independent days. Descriptive; body size was examined after the fact.}
\end{table}

\begin{table}[htbp]
\centering
\caption{Signal days in thirds by body-to-range ratio.}
\label{tab:auditthirds}
\begin{tabular}{@{}lrrrrrrrr@{}}
\toprule
 & \multicolumn{3}{c}{2013--2020} & \multicolumn{5}{c}{2021--2026} \\
\cmidrule(lr){2-4}\cmidrule(lr){5-9}
 & Ratio & Days & bps & Ratio & Days & bps & Trades & \$/trade \\
\midrule
Lowest third & 0.60--0.76 & 303 & $+3.95$ & 0.60--0.74 & 215 & $+1.84$ & 191 & $+19$ \\
Middle third & 0.76--0.93 & 302 & $+1.38$ & 0.74--0.87 & 214 & $+2.81$ & 201 & $+102$ \\
Highest third & 0.93--1.00 & 303 & $+1.91$ & 0.87--1.00 & 215 & $+2.44$ & 193 & $+80$ \\
\bottomrule
\end{tabular}
\tabnotes{Thirds are formed within each period. Index move: CFD, 15:51 open to
15:59 close, in the candle's direction, in basis points. Option result: dollars
per trade, \ITM{10} at the 15:51 ask, held to settlement. Descriptive.}
\end{table}

I looked at the candle's size, its body as a fraction of the index, only after
the fact, and it is in no plan. On signal
days a larger candle was followed by a larger move in its direction, in both
periods and in the option result, but not by the right direction more often, and
its rank correlations are smaller than its linear ones. A common volatility factor would produce this pattern, with larger candles in
more volatile minutes followed by larger gains and larger losses, and the
relation carried by a small number of large days. I treat it as a hypothesis
for the forward record.

The re-validation also found a difference between the signal feed and the
index. Every signal in this paper is read from the CFD candle, but a trader,
and the live alert, reads the SPX index candle. The two track each other closely in level and direction, but a candle's
body-to-range ratio depends on its shape, and the shapes differ. Real SPX
one-minute bars are available for 36 sessions in the autumn of 2026, and the CFD
series was extended over the same days. On those days the index candle showed a
strong body at 15:50 more often than the CFD candle (Table~\ref{tab:auditfeeds}).
Every CFD signal was also an index signal, in the same direction, so no
backtested signal was missed, but the index fired on six further days that the
backtests would not have traded. Across all one-minute candles from 15:30 to
15:59 the pattern is the same: an index candle in the 60--70\% range is a CFD
strong body less than half the time, and only above 80\% do the two nearly
always agree.

\begin{table}[htbp]
\centering
\caption{The SPX index candle against the CFD candle, autumn 2026.}
\label{tab:auditfeeds}
\begin{tabular}{@{}lrr@{\hspace{2.5em}}lrrr@{}}
\toprule
Index ratio & Candles & CFD $\ge 0.60$ & Cut-off & Days fired & CFD days & Others \\
\midrule
0.0--0.4 & 398 & 5\% & 0.60 & 20 & 14 & 6 \\
0.4--0.6 & 192 & 16\% & 0.65 & 18 & 14 & 4 \\
0.6--0.7 & 143 & 42\% & 0.70 & 17 & 14 & 3 \\
0.7--0.8 & 113 & 61\% & 0.75 & 16 & 14 & 2 \\
0.8--0.9 & 98 & 82\% & 0.80 & 16 & 14 & 2 \\
0.9--1.0 & 136 & 87\% & & & & \\
\bottomrule
\end{tabular}
\tabnotes{36 sessions, 2026-09-21 to 2026-11-09. Left: 1{,}080 one-minute candles from
15:30 to 15:59, by the index candle's body-to-range ratio, with the share that
were CFD strong bodies. Right: the 15:50 candle, by the cut-off applied to the
index candle; the CFD candle had 14 strong-body days, all among the index
candle's. The median 15:50 ratio was 0.65 on the index candle and 0.54 on the
CFD.}
\end{table}

A 60\% threshold applied to the index candle is therefore a looser rule than the
one that was tested; on this sample the tested rule fires about as often as a
threshold of roughly 70\% on the index candle. That correspondence concerns
frequency only, and the sample is a single autumn. Whether the additional days
pay was tested separately over 2021--2026 (Section~\ref{sec:realindex}): they did
not, and even above 70\% the two candles are not interchangeable.

The two ratios are also strongly correlated, at 0.90 for the 15:50 candle over
the 36 sessions and 0.76 across all 1{,}080 candles from 15:30 to 15:59. The least-squares line
through those candles, CFD ratio $= 0.10 + 0.72 \times$ index ratio, puts the
CFD's 0.60 at about 0.69 on the index. This is a second route to the
correspondence above, and it does not change it.

The comparison continues in a two-feed forward record. Since 9 November 2026 the live alert
records, every session, the 15:50 candle from both sources, each one's ratio and
band (below 0.60, 0.60--0.70, 0.70 or more) and the index move from the 15:51
open to the 15:59 close, and it labels each alerted signal with the index
candle's band. The record starts with the 36 sessions already on file, of which
only the 14 after 20 October lie outside every backtest
(Table~\ref{tab:auditforward}). It shows no difference yet: the index moved the
candle's way about as often on sessions without a signal as on sessions with
one, and on the six sessions where only the index candle reached 60\% it did so
four times. A correlation of the size found in the history needs on the order
of a thousand sessions to be detected, and the bands need a few hundred signals
before they can be compared. Section~\ref{sec:realindex} adds a more specific question for the record, whether sessions on which both candles show a strong body continue
to outperform sessions on which only the index candle does. Since the study of
Section~\ref{sec:spycandle}, the record also holds the SPY candle, and the alert
acts on the rule described at the end of that section.

\begin{table}[htbp]
\centering
\caption{The two-feed record, autumn 2026.}
\label{tab:auditforward}
\begin{tabular}{@{}lrrrrrr@{}}
\toprule
 & \multicolumn{3}{c}{Index candle} & \multicolumn{3}{c}{CFD candle} \\
\cmidrule(lr){2-4}\cmidrule(lr){5-7}
Body / range & Days & Right & Move & Days & Right & Move \\
\midrule
0.70 or more & 17 & 11 & $+1.53$ & 13 & 8 & $+1.55$ \\
0.60--0.70 & 3 & 2 & $+3.87$ & 1 & 1 & $+1.24$ \\
Below 0.60 & 16 & 10 & $+1.08$ & 22 & 14 & $+1.53$ \\
\bottomrule
\end{tabular}
\tabnotes{36 sessions, 2026-09-21 to 2026-11-09, by the band of the 15:50 candle
on each source; only the 14 sessions after 2026-10-20 lie outside every
backtest. ``Right'' counts sessions on which the index moved in the candle's
direction from the 15:51 open to the 15:59 close; ``Move'' is the mean of that
move, in index points. Descriptive; too few sessions to compare bands.}
\end{table}

% =====================================================================
\section{The rule on the index candle}
\label{sec:realindex}

The re-validation left one question open. Every signal in this paper was read
from the CFD candle, while a trader, and the live alert, reads the SPX index
candle, and Section~\ref{sec:audit} showed that the two differ in shape. A
further pre-registered study therefore read the 15:50 candle from an archive of
one-minute SPX index bars published by CNBC, the source of the live alert, and
traded it at the option prices of Section~\ref{sec:realprice}. The plan was
locked before any index candle had been joined with a price. The option
outcomes of these days were already known, so the study is, in the plan's words,
an outcome-informed, retrospective check of the signal's source: not an
independent confirmation and not a new holdout. Read from the index candle, the
rule did not meet its criterion, and the gain was confined to the days on which
both candles showed a strong body.

The archive holds a valid 15:50 bar on 1{,}310 of the
1{,}322 sessions from 2021-01-01 to 2026-10-20 that have option quotes and an
official close (99.1\%). Before the lock, and using price levels only, the
archived bar was compared with the broker's own index bar on the 22 sessions of
the period that have both, and all four prices were equal on 21. In every year
the archived 15:50 close lay nearer the CFD's 15:50 close than the CFD's closes a
minute earlier or later, which confirms the time stamps. The signal is the live
alert's test, computed in whole cents: a body of at least 60\% of the range,
split into a 60--70\% band and a band of 70\% or more. The trade is unchanged:
\ITM{10} in the candle's direction at the 15:51:00 ask, held to settlement. The
primary line asks whether the rule read from the index candle has a positive mean
per day with Newey--West $t \ge 2.0$. Three secondary lines, one for each band
and one for the difference between them per trade by Welch's test, also require
a Holm-adjusted $p \le 0.05$.

\begin{table}[htbp]
\centering
\small
\caption{The rule read from the SPX index candle, 2021--2026.}
\label{tab:realindex}
\begin{tabular}{@{}lrrrrrrrl@{}}
\toprule
Index candle & Trades & \$/trade & At mid & Won & \$/day & $t$ & Holm $p$ & Reading \\
\midrule
Body $\ge 60\%$ & 710 & $+\dollar39.44$ & $+\dollar78.14$ & 52\% & $+\dollar21.39$ & $+1.53$ &  & criterion not met \\
Body $\ge 70\%$ & 579 & $+\dollar53.73$ & $+\dollar91.20$ & 53\% & $+\dollar23.75$ & $+1.86$ & 0.094 & insufficient evidence \\
Body 60--70\% & 131 & $-\dollar23.69$ & $+\dollar20.44$ & 49\% & $-\dollar2.37$ & $-0.41$ & 0.685 & insufficient evidence \\
\bottomrule
\end{tabular}
\tabnotes{One \ITM{10} contract at the 15:51 ask, held to settlement, 2021-01-01
to 2026-10-20. The first row is the primary line. Per day is the mean of the
daily series (0 on days outside the cell); $t$ is Newey--West, lag 5. Holm $p$ is
adjusted across the three secondary lines, the third of which, the difference
between the bands, is given in the text. Pre-registered.}
\end{table}

Table~\ref{tab:realindex} gives the registered lines. Read
from the index candle, the rule made $+\dollar39.44$ a trade, against
$+\dollar67.24$ from the CFD candle on the same days (583 trades), and its daily
series fell short of the criterion ($t = 1.53$). The source-check criterion was
therefore not met. The band of 70\% or more came closer and also fell short. The
60--70\% band lost a little at the ask and made a little at the mid; under the
plan this is insufficient evidence in either direction, and the difference
between the two bands ($+\dollar77.41$ a trade, Welch $t = 1.20$) cannot be
distinguished from zero. By year (Table~\ref{tab:realindexyears}) the 60--70\%
band alternated between large losses and large gains with no pattern, the band of
70\% or more lost money only in 2023, and the rule as a whole made almost nothing
from 2022 to 2024.

\begin{table}[htbp]
\centering
\small
\caption{Where the two candles disagree, 2021--2026.}
\label{tab:realindexdisagree}
\begin{tabular}{@{}lrrrrrrrrr@{}}
\toprule
 & & & & & & \multicolumn{2}{c}{Index $\ge 70\%$} & \multicolumn{2}{c}{Index 60--70\%} \\
\cmidrule(lr){7-8}\cmidrule(lr){9-10}
Strong body on & Trades & Ask & Mid & Won & $t$ & Trades & Ask & Trades & Ask \\
\midrule
Both candles & 520 & $+\dollar80.95$ & $+\dollar120.14$ & 54\% & $+2.51$ & 481 & $+\dollar86$ & 39 & $+\dollar14$ \\
Index candle only & 190 & $-\dollar74.14$ & $-\dollar36.80$ & 45\% & $-1.68$ & 98 & $-\dollar107$ & 92 & $-\dollar40$ \\
CFD candle only & 63 & $-\dollar45.92$ & $-\dollar15.20$ & 51\% & $-0.73$ & 0 &  & 0 &  \\
\bottomrule
\end{tabular}
\tabnotes{\ITM{10} in the index candle's direction (in the CFD candle's direction
on days only the CFD candle signals), \dollar\ per trade at the ask and at the
mid; $t$ is a simple $t$-statistic. The two candles never signalled opposite
directions. Descriptive.}
\end{table}

The analysis of where the difference comes from is descriptive. The two candles
never pointed in opposite directions, but they often disagreed about
strength (Table~\ref{tab:realindexdisagree}), and the gain was confined to the
days on which both showed a strong body. The 190 trades taken on days when only
the index candle did lost $\dollar74.14$ a trade, and they include index candles
of 70\% or more. Reading the index candle therefore adds a set of days without
the effect, and these lower the mean. The correspondence of
Section~\ref{sec:audit} still holds as a statement about frequency: over 1{,}310
days the two ratios have a correlation of 0.78, and the CFD's 0.60 corresponds
to 0.69 on the index candle. It does not make the candles interchangeable, since
17\% of index candles of 70\% or more were below 60\% on the CFD. One
explanation, which I did not test, is that the CFD follows the futures, which
move first, while the index is computed from the last trades in 500 stocks and
can show a strong candle while it catches up.

By the size of the index candle's body (Table~\ref{tab:realindextenths}), the
option trade was profitable only from 80\% upward (438 trades, $+\dollar89$ a
trade); that cut was not planned. The index itself moved in the candle's
direction after candles of 70\% or more and not after candles of 60--70\%
(Table~\ref{tab:realindexmove}), so the directional effect is present on the index candle, and only the option
trade falls short.

\begin{table}[htbp]
\centering
\small
\caption{Option results by the index candle's body-to-range ratio, 2021--2026.}
\label{tab:realindextenths}
\begin{tabular}{@{}lrrrrr@{}}
\toprule
Body / range & Days & Trades & Ask & Mid & Won \\
\midrule
0.0--0.1 & 80 & 79 & $+\dollar78$ & $+\dollar118$ & 53\% \\
0.1--0.2 & 90 & 90 & $-\dollar21$ & $+\dollar17$ & 50\% \\
0.2--0.3 & 93 & 92 & $-\dollar23$ & $+\dollar17$ & 46\% \\
0.3--0.4 & 105 & 105 & $-\dollar42$ & $-\dollar5$ & 50\% \\
0.4--0.5 & 115 & 115 & $-\dollar53$ & $-\dollar17$ & 44\% \\
0.5--0.6 & 113 & 113 & $-\dollar12$ & $+\dollar28$ & 52\% \\
0.6--0.7 & 132 & 131 & $-\dollar24$ & $+\dollar20$ & 49\% \\
0.7--0.8 & 141 & 141 & $-\dollar56$ & $-\dollar15$ & 47\% \\
0.8--0.9 & 160 & 160 & $+\dollar67$ & $+\dollar104$ & 54\% \\
0.9--1.0 & 278 & 278 & $+\dollar102$ & $+\dollar137$ & 54\% \\
\bottomrule
\end{tabular}
\tabnotes{\ITM{10} in the candle's direction whatever the ratio, \dollar\ per
trade. A ratio of exactly 1 is in the last row. Descriptive.}
\end{table}

\begin{table}[htbp]
\centering
\small
\caption{The index move after the index candle, 2021--2026.}
\label{tab:realindexmove}
\begin{tabular}{@{}lrrrr@{}}
\toprule
Index candle & Days & bps & $t$ & Right \\
\midrule
Body $\ge 70\%$ & 640 & $+2.63$ & $+4.55$ & 57\% \\
Body 60--70\% & 141 & $+0.14$ & $+0.11$ & 51\% \\
Body $< 60\%$ & 650 & $+0.23$ & $+0.46$ & 53\% \\
\bottomrule
\end{tabular}
\tabnotes{15:51 open to 15:59 close of the archived index bars, in the candle's
direction, in basis points; $t$ is Newey--West, lag 5; ``Right'' is the share of
non-zero moves in the candle's direction. Descriptive.}
\end{table}

\begin{table}[htbp]
\centering
\renewcommand{\tabfont}{\fontsize{9}{11}\selectfont}
\small
\caption{The rule read from the index candle, by year.}
\label{tab:realindexyears}
\begin{tabular}{@{}lrrrrrr@{}}
\toprule
Signal on the index candle & 2021 & 2022 & 2023 & 2024 & 2025 & 2026 \\
\midrule
Body $\ge 60\%$ & $+\dollar105$ (88) & $+\dollar7$ (112) & $+\dollar5$ (136) & $+\dollar7$ (125) & $+\dollar51$ (134) & $+\dollar83$ (115) \\
Body $\ge 70\%$ & $+\dollar218$ (63) & $+\dollar45$ (91) & $-\dollar19$ (110) & $+\dollar32$ (102) & $+\dollar19$ (118) & $+\dollar104$ (95) \\
Body 60--70\% & $-\dollar179$ (25) & $-\dollar157$ (21) & $+\dollar109$ (26) & $-\dollar104$ (23) & $+\dollar285$ (16) & $-\dollar16$ (20) \\
\bottomrule
\end{tabular}
\tabnotes{\dollar\ per trade at the ask, with the number of trades in
parentheses. Descriptive.}
\end{table}

The study has limits. The days and their option outcomes are those of the
earlier confirmation; only the candle changed, so the registered lines test the source of
the signal and nothing else. The archive is a vendor's stored bar, which need not
equal the bar shown live, and 12 sessions had none; on those the CFD rule's known
result was ordinary ($+\dollar22$ a day against $+\dollar30$ on the others). The
disagreement table, the tenths and the 80\% cut were read after the result. As in
every real-price study, entry is at the quote of the instant the candle closes.

% =====================================================================
\section{The rule on the SPY candle}
\label{sec:spycandle}

The index-candle study left a practical problem. The candle that met its criterion,
the CFD's, cannot be seen live, and the candle that can, the index's, did not
meet it. The CFD follows the futures, and SPY, the exchange-traded fund, trades
live and follows them as well. A further pre-registered study therefore read the
15:50 candle from archived one-minute SPY bars and traded the same SPXW option at
the same prices. SPY was the third candle source evaluated against the same
option outcomes, and the study's second line was chosen after the agreement of
two candles had been seen to do well. The plan was locked before any SPY candle
was joined with a price, which prevents further adaptation but does not remove
that selection. Its lines are therefore retrospective numerical criteria, not an
independent confirmation, a new holdout or an established trading edge. Read from
the SPY candle, the rule met those criteria, and the gain was again confined to
sessions on which all three candles showed a strong body.

The archive, from CNBC as for the index, holds a valid SPY
15:50 bar on 1{,}312 of the 1{,}322 sessions with option quotes and an official
close (99.2\%). Checks made before the lock, on price levels only, showed that
the bar labelled 15:50 is the same minute as the index bar in every year, that
the bars are traded prices and not dividend-adjusted (their 15:59 close equals
the unadjusted daily close at the median in every year), and that 16\% of the
prices are sub-penny prints, which are rounded to the cent. They also showed
that an ETF candle depends on its vendor: against a second vendor's SPY bars on
22 sessions, the four prices of the 15:50 bar agreed to the cent on 14, with
differences of up to \$0.02 elsewhere. The signal is the same whole-cent test
applied to the SPY bar, and the trade is \ITM{10} in the SPY candle's direction
at the 15:51:00 ask, held to settlement. The primary line asks whether the daily
series has a positive mean with Newey--West $t \ge 2.0$. Two secondary lines, the
SPY and index candles both signalling in the same direction and the SPY candle's
band of 70\% or more, also require a Holm-adjusted $p \le 0.05$.

\begin{table}[htbp]
\centering
\renewcommand{\tabfont}{\fontsize{9}{11}\selectfont}
\small
\caption{The rule read from the SPY candle, 2021--2026.}
\label{tab:spycandle}
\begin{tabular}{@{}lrrrrrrrl@{}}
\toprule
Signal & Trades & \$/trade & At mid & Won & \$/day & $t$ & Holm $p$ & Reading \\
\midrule
SPY body $\ge 60\%$ & 523 & $+\dollar71.74$ & $+\dollar108.51$ & 53\% & $+\dollar28.62$ & $+2.35$ &  & criterion met \\
SPY and index both $\ge 60\%$ & 504 & $+\dollar80.44$ & $+\dollar117.34$ & 54\% & $+\dollar31.21$ & $+2.56$ & 0.0053 & criterion met \\
SPY body $\ge 70\%$ & 369 & $+\dollar103.40$ & $+\dollar137.67$ & 55\% & $+\dollar29.08$ & $+2.84$ & 0.0045 & criterion met \\
\midrule
SPY body 60--70\% & 154 & $-\dollar4.14$ & $+\dollar38.66$ & 49\% & & & & descriptive \\
CFD rule, same days & 582 & $+\dollar68.84$ & $+\dollar107.33$ & 54\% & & & & descriptive \\
\bottomrule
\end{tabular}
\tabnotes{One \ITM{10} contract at the 15:51 ask in the SPY candle's direction,
held to settlement, 2021-01-01 to 2026-10-20. The first row is the primary line
and the next two are the secondary lines; Holm $p$ is adjusted across those two.
Per day is the mean of the daily series; $t$ is Newey--West, lag 5. The last two
rows are descriptive. Pre-registered, retrospective.}
\end{table}

Table~\ref{tab:spycandle} shows that all three registered lines met their
retrospective criteria. Read from the SPY candle, the rule made
$+\dollar71.74$ a trade ($t = 2.35$), about what the CFD candle made on the same
days and nearly twice what the index candle made. The band of 70\% or more was
the strongest line and was positive in each of the six years
(Table~\ref{tab:spyyears}); the 60--70\% band made nothing, as on the other two
sources.

\begin{table}[htbp]
\centering
\renewcommand{\tabfont}{\fontsize{9}{11}\selectfont}
\small
\caption{The rule read from the SPY candle, by year.}
\label{tab:spyyears}
\begin{tabular}{@{}lrrrrrr@{}}
\toprule
Signal & 2021 & 2022 & 2023 & 2024 & 2025 & 2026 \\
\midrule
SPY body $\ge 60\%$ & $+\dollar201$ (51) & $+\dollar46$ (78) & $+\dollar47$ (98) & $+\dollar60$ (98) & $+\dollar44$ (107) & $+\dollar93$ (91) \\
SPY and index both $\ge 60\%$ & $+\dollar210$ (49) & $+\dollar81$ (75) & $+\dollar50$ (92) & $+\dollar71$ (95) & $+\dollar44$ (103) & $+\dollar92$ (90) \\
SPY body $\ge 70\%$ & $+\dollar281$ (34) & $+\dollar141$ (46) & $+\dollar30$ (70) & $+\dollar94$ (76) & $+\dollar95$ (79) & $+\dollar84$ (64) \\
SPY body 60--70\% & $+\dollar40$ (17) & $-\dollar91$ (32) & $+\dollar91$ (28) & $-\dollar56$ (22) & $-\dollar102$ (28) & $+\dollar117$ (27) \\
\bottomrule
\end{tabular}
\tabnotes{\dollar\ per trade at the ask, with the number of trades in
parentheses. Descriptive.}
\end{table}

The comparison of the three candles is descriptive. A strong SPY candle is
nearly always a strong CFD candle: 94\% of SPY candles of 70\% or more were CFD
signals (Table~\ref{tab:spycross}). The gain was confined to sessions on which
all three candles showed a strong body (435 trades, $+\dollar112.78$ a trade;
Table~\ref{tab:spythree}), and every other combination lost on average,
including the index and CFD candles signalling without SPY. No single row of
that table is significant on its own, but together they repeat the finding of Section~\ref{sec:realindex} that a
strong body on only some of the candles carries no edge. By the size of the SPY candle's body (Table~\ref{tab:spytenths}), the
trade was profitable in every tenth from 70\% upward, a lower cut than the index
candle's 80\%; below 70\% only one tenth, 0.1--0.2, was positive. The index moved
in the SPY candle's direction after SPY candles of 70\% or more
(Table~\ref{tab:spymove}).

\begin{table}[htbp]
\centering
\small
\caption{Where a SPY candle's band falls on the other two candles.}
\label{tab:spycross}
\begin{tabular}{@{}llrrrr@{}}
\toprule
Other candle & SPY band & Days & Other $< 60\%$ & Other 60--70\% & Other $\ge 70\%$ \\
\midrule
CFD & SPY $\ge 70\%$ & 369 & 6\% & 11\% & 83\% \\
CFD & SPY 60--70\% & 155 & 35\% & 33\% & 32\% \\
CFD & SPY $< 60\%$ & 788 & 83\% & 7\% & 10\% \\
\midrule
Index & SPY $\ge 70\%$ & 366 & 0\% & 2\% & 98\% \\
Index & SPY 60--70\% & 154 & 9\% & 18\% & 73\% \\
Index & SPY $< 60\%$ & 780 & 74\% & 12\% & 14\% \\
\bottomrule
\end{tabular}
\tabnotes{The share of each SPY band's sessions in each band of the CFD candle
and of the index candle, 2021--2026. Descriptive.}
\end{table}

\begin{table}[htbp]
\centering
\small
\caption{Which of the three candles show a strong body, 2021--2026.}
\label{tab:spythree}
\begin{tabular}{@{}cccrrrrrr@{}}
\toprule
SPY & Index & CFD & Days & Trades & Ask & Mid & Won & $t$ \\
\midrule
yes & yes & yes & 436 & 435 & $+\dollar112.78$ & $+\dollar150.18$ & 56\% & $+3.11$ \\
yes & yes & -- & 69 & 69 & $-\dollar123.42$ & $-\dollar89.66$ & 36\% & $-1.73$ \\
yes & -- & yes & 9 & 9 & $-\dollar118.42$ & $-\dollar91.20$ & 33\% & $-0.80$ \\
yes & -- & -- & 6 & 6 & $-\dollar401.42$ & $-\dollar378.09$ & 33\% & $-1.25$ \\
-- & yes & yes & 82 & 82 & $-\dollar77.58$ & $-\dollar28.55$ & 44\% & $-1.20$ \\
-- & yes & -- & 121 & 121 & $-\dollar46.04$ & $-\dollar6.66$ & 50\% & $-0.82$ \\
-- & -- & yes & 54 & 54 & $-\dollar33.83$ & $-\dollar2.53$ & 54\% & $-0.49$ \\
-- & -- & -- & 523 & 0 & & & & \\
\bottomrule
\end{tabular}
\tabnotes{\ITM{10} in the direction of the first signalling candle in the order
SPY, index, CFD; 1{,}300 sessions with all three candles; \dollar\ per trade at
the ask and at the mid; $t$ is a simple $t$-statistic. No session had signalling
candles in opposite directions. Descriptive.}
\end{table}

\begin{table}[htbp]
\centering
\small
\caption{Option results by the SPY candle's body-to-range ratio, 2021--2026.}
\label{tab:spytenths}
\begin{tabular}{@{}lrrrrr@{}}
\toprule
Body / range & Days & Trades & Ask & Mid & Won \\
\midrule
0.0--0.1 & 91 & 91 & $-\dollar51$ & $-\dollar5$ & 53\% \\
0.1--0.2 & 110 & 109 & $+\dollar119$ & $+\dollar161$ & 61\% \\
0.2--0.3 & 124 & 124 & $-\dollar28$ & $+\dollar10$ & 45\% \\
0.3--0.4 & 144 & 143 & $-\dollar87$ & $-\dollar49$ & 45\% \\
0.4--0.5 & 137 & 137 & $-\dollar37$ & $+\dollar5$ & 53\% \\
0.5--0.6 & 161 & 161 & $-\dollar74$ & $-\dollar35$ & 45\% \\
0.6--0.7 & 155 & 154 & $-\dollar4$ & $+\dollar39$ & 49\% \\
0.7--0.8 & 145 & 145 & $+\dollar90$ & $+\dollar124$ & 54\% \\
0.8--0.9 & 125 & 125 & $+\dollar137$ & $+\dollar174$ & 56\% \\
0.9--1.0 & 99 & 99 & $+\dollar81$ & $+\dollar112$ & 55\% \\
\bottomrule
\end{tabular}
\tabnotes{\ITM{10} in the candle's direction whatever the ratio, \dollar\ per
trade. A ratio of exactly 1 is in the last row. Descriptive.}
\end{table}

\begin{table}[htbp]
\centering
\small
\caption{The index move after the SPY candle, 2021--2026.}
\label{tab:spymove}
\begin{tabular}{@{}lrrrr@{}}
\toprule
SPY candle & Days & bps & $t$ & Right \\
\midrule
SPY body $\ge 70\%$ & 399 & $+3.42$ & $+4.70$ & 58\% \\
SPY body 60--70\% & 166 & $+1.26$ & $+1.20$ & 56\% \\
SPY body $< 60\%$ & 835 & $+0.35$ & $+0.74$ & 55\% \\
\bottomrule
\end{tabular}
\tabnotes{15:51 open to 15:59 close of the archived index bars, in the SPY
candle's direction, in basis points; $t$ is Newey--West, lag 5; ``Right'' is the
share of non-zero moves in the candle's direction. Descriptive.}
\end{table}

Several limits apply. SPY is the third source tried on the same outcomes, so its
$t$-statistics overstate what a fresh sample would support, and the days are not
new. The result belongs to one vendor's archived SPY candle: a second vendor's
candle differed by a cent or two on about a third of the sessions compared,
enough to move a borderline candle across a band edge. The 10 sessions without
an archived SPY bar were poor days for the CFD rule ($-\dollar60$ a day against
$+\dollar31$ on the rest), so the estimates are slightly flattered. The
three-candle table, the tenths and the yearly figures were read after the
result, and entry is at the quote of the instant the candle closes. Since the
run, the live alert reads the SPY candle as well as the index candle and records
both, so the comparison continues on new days.

The rule now in use is descriptive and was not registered. After the SPY study
the live alert was changed to act on one rule
built from its lines: the SPY 15:50 candle in its band of 70\% or more, and the index candle at 60\% or more in the
same direction; \ITM{10} in that direction at the 15:51 ask, held to settlement. This is the intersection of two
cells the study reported, not a registered line of its own, and it is nearly the 70\% band itself: of that band's
369 trades, 3 had no index bar and 1 had an index candle below
60\%, which leaves 365. On 2021-01-01 to 2026-10-20 those trades made $+\dollar101.33$ each and
$+\dollar36{,}986$ in all, with 54.5\% profitable, an average gain of \dollar593 against an
average loss of \dollar488, a worst trade of $-\dollar2{,}210$, a longest run of
7 losing trades and a worst peak-to-trough decline of \dollar5{,}271
(Table~\ref{tab:rulenow}). The rule traded on about 30\%
of sessions once expiries became daily. These figures describe days that helped to shape the rule, with the entry
at the instant the candle closes, and they are not evenly spread: from 2026-08-10 to 2026-10-20, the last ten weeks
of the sample, the rule lost \dollar974 over 18 trades, 7 of them profitable. The
7 signal days after the sample (2026-10-22 to 2026-11-06) were too few to mean anything: 4 priced from the collector's one-minute broker bars made $+\dollar1{,}214$ and 3 priced from OPRA asks bought afterwards made $+\dollar562$, $+\dollar1{,}775$ in all with 5 of 7 profitable. The alert now sends two messages a session, whether the candles look like arming at about 15:50:40
and the decision right after 15:51:00, and logs every session's three candles and the outcome.

\begin{table}[htbp]
\centering
\small
\caption{The rule now in use, by year.}
\label{tab:rulenow}
\begin{tabular}{@{}lrrrrrr@{}}
\toprule
Year & Sessions & Trades & Share & \$/trade & Total & Won \\
\midrule
2021 & 251 & 33 & 13\% & $+\dollar291$ & $+\dollar9{,}595$ & 20 \\
2022 & 250 & 45 & 18\% & $+\dollar140$ & $+\dollar6{,}309$ & 28 \\
2023 & 248 & 68 & 27\% & $+\dollar15$ & $+\dollar1{,}050$ & 36 \\
2024 & 249 & 76 & 31\% & $+\dollar94$ & $+\dollar7{,}150$ & 40 \\
2025 & 247 & 79 & 32\% & $+\dollar95$ & $+\dollar7{,}527$ & 45 \\
2026 & 201 & 64 & 32\% & $+\dollar84$ & $+\dollar5{,}355$ & 30 \\
\midrule
All & 1{,}446 & 365 & 25\% & $+\dollar101$ & $+\dollar36{{,}}986$ & 199 \\
\bottomrule
\end{tabular}
\tabnotes{SPY 15:50 candle of 70\% or more and index candle of 60\% or more in the
same direction, \ITM{10} at the 15:51 ask, held to settlement, one contract.
Sessions are full sessions; expiries were not daily before 2022. Won is the number
of profitable trades. Descriptive; not a registered line.}
\end{table}

% =====================================================================
\section{A search over closing candles}
\label{sec:search}

The studies so far fixed the candle and its cut and asked whether that rule
worked. A last pre-registered study asked a broader question, namely which setting
among the closing candles, cuts, directions and strikes near the rule does best
on real asks, and whether the best one is better than luck. The plan was locked
after every earlier study, and two of its three candle sources, the SPY candle
and the agreement of the SPY and index candles, were chosen after their results
had been seen. It therefore guards against further selection but not against
that history, and it set aside in advance every 15:50 setting, because those had
already been examined on the later years and could not be tested on them afresh.
The search chose a 15:50 setting again, so nothing went to the test, and no other
closing minute beat the 15:50 candle on the tuning years.

The search covers 7{,}128 settings: three candle
sources (the index candle, the SPY candle, and the two agreeing in direction),
the eleven one-minute candles from 15:45 to 15:55, twelve body cuts from 40\% to
95\%, trading with or against the candle, and nine strikes from \ITM{20} to
\OTM{20}. Each trade buys one SPXW contract at the ask in the snapshot that
opens the next minute and holds it to settlement. A setting is scored on
2021--2023 by its mean profit per day averaged with the settings one cut above
and below, so that an isolated spike does not win. To be eligible it needs at
least 100 trades, at least 40 in each half of the period, a positive mean in each
half, and prices for at least 95\% of its signals; 387 settings
qualified. The best one would be tested once on 2024--2026 only if it was not a
15:50 setting, if a luck diagnostic gave $p \le 0.05$, and if it beat the
index-candle rule of Section~\ref{sec:realindex} on the tuning days. The luck
diagnostic repeats the whole search 195 times with every candle signal
moved by a whole number of sessions against fixed outcomes, which keeps each
signal's own pattern but breaks its link to the outcome. Before any statistic,
the code reproduced the trades of the two candle studies exactly
(336 and 227 on the tuning days, 374
and 296 on the test days). It also passed synthetic self-tests and a full trial on synthetic data.

\begin{table}[htbp]
\centering
\small
\caption{The best settings on 2021--2023.}
\label{tab:searchtune}
\begin{tabular}{@{}llllrrrrrr@{}}
\toprule
Source & Candle & Body & Side & Strike & Score & \$/day & Trades & First half & Second half \\
\midrule
Index & 15:50 & $\ge 80\%$ & With & \ITM{20} & $+\dollar37.46$ & $+\dollar39.42$ & 193 & $+\dollar77.15$ & $+\dollar1.08$ \\
Index & 15:50 & $\ge 80\%$ & With & \ITM{15} & $+\dollar33.87$ & $+\dollar35.50$ & 192 & $+\dollar69.96$ & $+\dollar0.59$ \\
Index & 15:54 & $\ge 60\%$ & With & \ITM{5} & $+\dollar33.55$ & $+\dollar32.43$ & 299 & $+\dollar50.12$ & $+\dollar14.33$ \\
Both & 15:54 & $\ge 40\%$ & With & \ITM{5} & $+\dollar32.88$ & $+\dollar38.17$ & 325 & $+\dollar53.11$ & $+\dollar23.03$ \\
Index & 15:54 & $\ge 65\%$ & With & \ITM{5} & $+\dollar32.47$ & $+\dollar32.71$ & 261 & $+\dollar50.36$ & $+\dollar14.71$ \\
Both & 15:50 & $\ge 65\%$ & With & \ITM{20} & $+\dollar32.19$ & $+\dollar34.72$ & 179 & $+\dollar67.80$ & $+\dollar1.09$ \\
Index & 15:54 & $\ge 70\%$ & With & \ITM{5} & $+\dollar31.94$ & $+\dollar32.26$ & 224 & $+\dollar45.21$ & $+\dollar18.97$ \\
Index & 15:50 & $\ge 80\%$ & With & \ITM{5} & $+\dollar31.88$ & $+\dollar33.85$ & 192 & $+\dollar67.00$ & $+\dollar0.27$ \\
SPY & 15:54 & $\ge 40\%$ & With & \ITM{5} & $+\dollar31.46$ & $+\dollar35.46$ & 342 & $+\dollar49.97$ & $+\dollar20.76$ \\
Both & 15:54 & $\ge 45\%$ & With & \ITM{5} & $+\dollar31.40$ & $+\dollar27.58$ & 283 & $+\dollar45.98$ & $+\dollar8.75$ \\
\bottomrule
\end{tabular}
\tabnotes{The ten eligible settings with the highest score, of 387
eligible among 7{,}128. Score is the mean profit per day averaged with the
same setting one cut above and below; \$/day is the setting's own mean over
625 sessions. The halves are 2021-01-01 to 2022-09-30 and 2022-10-01 to
2023-12-31. One contract at the ask, held to settlement, \$0.09 in fees.
Pre-registered.}
\end{table}

Table~\ref{tab:searchtune} shows that the search chose a 15:50 setting: the index
candle with a body of 80\% or more, traded in its direction with \ITM{20}. Because 15:50 settings were set aside, nothing
went to the test, and this choice can be checked only on later sessions. It also
missed the luck bar: the shifted searches reached its score in 9 of
195 cases ($p = 0.051$; median best score $+\dollar20.11$, 95th
percentile $+\dollar37.38$). Its gain came almost entirely in the first half
of the tuning period ($+\dollar77.15$ a day against $+\dollar1.08$).
By candle minute (Table~\ref{tab:searchminutes}), 15:50 led and 15:54 came
second, every other minute scored below both, and two minutes had no eligible
setting at all.

\begin{table}[htbp]
\centering
\small
\caption{The best eligible setting for each candle, 2021--2023.}
\label{tab:searchminutes}
\begin{tabular}{@{}llllrrrr@{}}
\toprule
Candle & Source & Body & Side & Strike & Score & \$/day & Trades \\
\midrule
15:45 & Index & $\ge 80\%$ & With & \ITM{20} & $+\dollar21.46$ & $+\dollar27.65$ & 166 \\
15:46 & Index & $\ge 40\%$ & With & \ITM{5} & $+\dollar1.48$ & $+\dollar6.81$ & 430 \\
15:47 & Index & $\ge 70\%$ & With & \ITM{5} & $-\dollar0.19$ & $+\dollar4.25$ & 235 \\
15:48 & Index & $\ge 40\%$ & With & \ATM & $+\dollar19.30$ & $+\dollar19.41$ & 416 \\
15:49 & Index & $\ge 90\%$ & With & \ATM & $+\dollar5.22$ & $+\dollar7.75$ & 104 \\
15:50 & Index & $\ge 80\%$ & With & \ITM{20} & $+\dollar37.46$ & $+\dollar39.42$ & 193 \\
15:51 & \multicolumn{4}{l}{No eligible setting} & & & \\
15:52 & Both & $\ge 50\%$ & Against & \ITM{5} & $+\dollar12.00$ & $+\dollar13.31$ & 237 \\
15:53 & Index & $\ge 80\%$ & With & \ATM & $+\dollar23.40$ & $+\dollar23.69$ & 177 \\
15:54 & Index & $\ge 60\%$ & With & \ITM{5} & $+\dollar33.55$ & $+\dollar32.43$ & 299 \\
15:55 & \multicolumn{4}{l}{No eligible setting} & & & \\
\bottomrule
\end{tabular}
\tabnotes{As in Table~\ref{tab:searchtune}, the highest-scoring eligible
setting among those that read the given candle. Entry is at the next minute.
Pre-registered.}
\end{table}

The ranking over the whole period is descriptive. After the decision,
the grid was ranked on all of 2021--2026 by mean profit per day, among settings
with at least 200 trades and prices for 95\% of their signals. All fifteen
leading settings were 15:50 candles traded in their direction, 10 of them
with \ITM{20} (Table~\ref{tab:searchwhole}). The first, Both 15:50, $\ge 60\%$, with, \ITM{20},
made $+\dollar34.74$ a day ($+\dollar89.44$ a trade,
$t = 2.71$, 505 trades), above the 95th percentile of the
same ranking repeated on 427 shifted signals
($+\dollar30.67$). The ranking reuses days already seen and is not a
test. Table~\ref{tab:searchleads} follows the 15:50 settings examined in
earlier sections across the two periods. Every one made money in both, the
60\% rule least on the tuning days, and on 2024--2026 only the SPY candle at
70\% with \ITM{20} reached $t = 2$. Paying the ask was costly throughout: at
the mid, the search's choice made $+\dollar61.03$ a day on
2021--2023 and $+\dollar41.32$ on 2024--2026, against
$+\dollar39.42$ and $+\dollar29.58$ at the
ask.

\begin{table}[htbp]
\centering
\small
\caption{The leading settings over 2021--2026, all on the 15:50 candle.}
\label{tab:searchwhole}
\begin{tabular}{@{}lllrrrrr@{}}
\toprule
Source & Body & Strike & \$/day & $t$ & Trades & \$/trade & Won \\
\midrule
Both & $\ge 60\%$ & \ITM{20} & $+\dollar34.74$ & $+2.71$ & 505 & $+\dollar89.44$ & 55\% \\
Index & $\ge 85\%$ & \ITM{20} & $+\dollar34.23$ & $+3.23$ & 365 & $+\dollar122.87$ & 55\% \\
Index & $\ge 80\%$ & \ITM{20} & $+\dollar34.20$ & $+2.92$ & 438 & $+\dollar102.29$ & 55\% \\
Both & $\ge 65\%$ & \ITM{20} & $+\dollar33.97$ & $+2.85$ & 429 & $+\dollar102.93$ & 55\% \\
SPY & $\ge 65\%$ & \ITM{20} & $+\dollar33.44$ & $+2.83$ & 446 & $+\dollar98.37$ & 55\% \\
SPY & $\ge 70\%$ & \ITM{20} & $+\dollar33.08$ & $+3.11$ & 369 & $+\dollar117.61$ & 55\% \\
Index & $\ge 75\%$ & \ITM{20} & $+\dollar32.23$ & $+2.67$ & 500 & $+\dollar84.44$ & 54\% \\
Both & $\ge 70\%$ & \ITM{20} & $+\dollar32.20$ & $+3.01$ & 357 & $+\dollar117.25$ & 55\% \\
Both & $\ge 60\%$ & \ITM{15} & $+\dollar32.10$ & $+2.53$ & 503 & $+\dollar82.84$ & 55\% \\
SPY & $\ge 60\%$ & \ITM{20} & $+\dollar31.92$ & $+2.49$ & 524 & $+\dollar79.92$ & 54\% \\
\bottomrule
\end{tabular}
\tabnotes{The ten settings with the highest mean profit per day on 2021-01-01 to
2026-10-20 among 3{,}887 with at least 200 trades and prices
for 95\% of their signals; all trade in the candle's direction. $t$ is
Newey--West, lag 5. Ranked after the tuning decision on days already seen.
Descriptive.}
\end{table}

\begin{table}[htbp]
\centering
\small
\caption{The 15:50 settings examined earlier, by period.}
\label{tab:searchleads}
\begin{tabular}{@{}lrrrrrr@{}}
\toprule
 & \multicolumn{3}{c}{2021--2023} & \multicolumn{3}{c}{2024--2026} \\
\cmidrule(lr){2-4}\cmidrule(l){5-7}
Setting & \$/day & $t$ & Trades & \$/day & $t$ & Trades \\
\midrule
Index $\ge 60\%$, \ITM{10} (Section~\ref{sec:realindex}) & $+\dollar17.54$ & $+0.87$ & 336 & $+\dollar24.80$ & $+1.27$ & 374 \\
\addlinespace
Index $\ge 80\%$, \ITM{10} & $+\dollar37.26$ & $+2.49$ & 193 & $+\dollar23.05$ & $+1.39$ & 245 \\
Index $\ge 80\%$, \ITM{20} & $+\dollar39.42$ & $+2.50$ & 193 & $+\dollar29.58$ & $+1.73$ & 245 \\
Index $\ge 85\%$, \ITM{10} & $+\dollar35.50$ & $+2.73$ & 161 & $+\dollar22.85$ & $+1.47$ & 204 \\
Index $\ge 85\%$, \ITM{20} & $+\dollar38.29$ & $+2.83$ & 161 & $+\dollar30.64$ & $+1.92$ & 204 \\
SPY $\ge 70\%$, \ITM{10} & $+\dollar29.28$ & $+2.29$ & 150 & $+\dollar28.91$ & $+1.85$ & 219 \\
SPY $\ge 70\%$, \ITM{20} & $+\dollar30.56$ & $+2.26$ & 150 & $+\dollar35.32$ & $+2.20$ & 219 \\
SPY $\ge 80\%$, \ITM{10} & $+\dollar22.50$ & $+2.12$ & 85 & $+\dollar16.22$ & $+1.22$ & 139 \\
SPY $\ge 80\%$, \ITM{20} & $+\dollar23.31$ & $+2.13$ & 85 & $+\dollar21.11$ & $+1.53$ & 139 \\
Both $\ge 60\%$, \ITM{10} & $+\dollar34.52$ & $+2.11$ & 216 & $+\dollar28.30$ & $+1.58$ & 288 \\
Both $\ge 60\%$, \ITM{20} & $+\dollar32.54$ & $+1.84$ & 217 & $+\dollar36.69$ & $+1.99$ & 288 \\
\bottomrule
\end{tabular}
\tabnotes{The 15:50 candle traded in its direction, one contract at the ask,
held to settlement. The first row is the index-candle rule; the rest are the
settings the plan set aside as examined. ``Both'' requires the index and SPY
candles to pass the cut in the same direction. $t$ is Newey--West, lag 5.
Descriptive.}
\end{table}

The question of which candle predicts the index move is also descriptive.
Without options, the same signals can be scored on the index move itself, from the next
minute's open to the 15:59 close in the candle's direction
(Table~\ref{tab:searchmove}). On 2021--2026 the ten largest $t$-statistics all
belonged to 15:50 candles, led by the index candle at 85\%
($+3.52$ basis points over 408 sessions, $t = 5.21$).
On 2019--2020, with the index candle alone, no 15:50 setting was among the ten
largest, and the largest $t$ (2.90, on 23 sessions) was below
the median of its shifted benchmark. The registered benchmark for this table,
the 95th percentile of the largest $t$ on shifted signals (5.57
on 2021--2026), was not cleared. A check chosen after the run shows why the bar
is high: in 62\% of the shifted
tables the largest $t$ came from a setting with fewer than 100 sessions, and
among settings with at least 300 sessions the shifted 95th
percentile was 3.40. The 15:50 result is
therefore strong for its sample size, but it did not pass the bar fixed in
advance.

\begin{table}[htbp]
\centering
\small
\caption{The index move after each closing candle.}
\label{tab:searchmove}
\begin{tabular}{@{}lllrrrr@{}}
\toprule
Source & Candle & Body & Sessions & bps & $t$ & Right \\
\midrule
\multicolumn{7}{@{}l}{\emph{Panel A: 2019--2020, index candle only}} \\
Index & 15:47 & $\ge 95\%$ & 23 & $+6.78$ & $+2.90$ & 57\% \\
Index & 15:53 & $\ge 90\%$ & 31 & $+7.95$ & $+2.85$ & 61\% \\
Index & 15:48 & $\ge 75\%$ & 60 & $+5.11$ & $+2.78$ & 60\% \\
Index & 15:48 & $\ge 80\%$ & 49 & $+5.01$ & $+2.65$ & 59\% \\
Index & 15:48 & $\ge 70\%$ & 73 & $+4.32$ & $+2.57$ & 59\% \\
\addlinespace
\multicolumn{7}{@{}l}{\emph{Panel B: 2021--2026, all three sources}} \\
Index & 15:50 & $\ge 85\%$ & 408 & $+3.52$ & $+5.21$ & 58\% \\
Index & 15:50 & $\ge 80\%$ & 490 & $+3.35$ & $+5.21$ & 59\% \\
Index & 15:50 & $\ge 75\%$ & 555 & $+2.99$ & $+5.01$ & 58\% \\
Both & 15:50 & $\ge 65\%$ & 468 & $+3.38$ & $+4.94$ & 59\% \\
SPY & 15:50 & $\ge 65\%$ & 482 & $+3.26$ & $+4.91$ & 59\% \\
Both & 15:50 & $\ge 60\%$ & 550 & $+2.97$ & $+4.79$ & 58\% \\
SPY & 15:50 & $\ge 70\%$ & 399 & $+3.42$ & $+4.70$ & 58\% \\
SPY & 15:50 & $\ge 60\%$ & 565 & $+2.78$ & $+4.61$ & 57\% \\
Both & 15:50 & $\ge 70\%$ & 390 & $+3.43$ & $+4.59$ & 58\% \\
Both & 15:50 & $\ge 55\%$ & 629 & $+2.61$ & $+4.58$ & 57\% \\
\bottomrule
\end{tabular}
\tabnotes{The settings with the largest $t$, of 132 on 2019--2020 and
396 on 2021--2026: the index move from the open of the minute after the
candle to the 15:59 close, in the candle's direction, in basis points, on
sessions on which the candle passed the cut. $t$ is Newey--West, lag 5;
``Right'' is the share of non-zero moves in the candle's direction. Shifted
benchmark for the largest $t$: median 3.26 and 95th percentile
6.39 on 2019--2020, 3.00 and
5.57 on 2021--2026. Descriptive.}
\end{table}

The search also has limits. It reuses the 2021--2026 outcomes of the earlier
studies, and two of its candle sources were chosen with knowledge of them, so
its luck diagnostic corrects for the 7{,}128 settings but not for that
history. Its choice was not clearly beyond luck and earned almost all of its
tuning gain before October 2022. The search is still useful as a comparison. No other closing minute, cut,
direction or strike beat the 15:50 candle on 2021--2023, and settings that ask
for a cleaner candle (80\% or more on the index, 70\% or more on SPY, or the
two candles agreeing) made more than the 60\% rule on 2021--2023 and between
$+\dollar16.22$ and $+\dollar36.69$ a day on 2024--2026, against its
$+\dollar24.80$. Those are hypotheses that only sessions after
November 2026 can test.

% Generated by backtests/spx_candle_imbalance/make_tex.py (spx_candle_1550/imbalance_section_draft.tex); do not edit numbers by hand.
% =====================================================================
\section{Closing-auction imbalances}
\label{sec:imbalance}



From 15:50~ET the New York Stock Exchange publishes the closing-auction imbalance of the stocks it lists, and Nasdaq
disseminates an early imbalance indicator for its own stocks every ten seconds.\footnote{The NYSE moved its
market-on-close cut-off and imbalance publication from 15:45 to 15:50 on 1 April 2019 (SEC Release No.~34-85021).
Nasdaq's early indicator from 15:50 was scheduled for the fourth quarter of 2019 (Release No.~34-87179). Samples of
the exchanges' messages on 30 December 2020 and 6 November 2026, outside the sample, confirmed both schedules.} The
15:50 candle is therefore the first full minute in which the market can react to the published imbalance. That
suggests an explanation of the signal: traders act on the publication during the 15:50 minute, and the index keeps
drifting the same way while the auction absorbs the imbalance. If that channel mattered, the rule should be right
more often on days when its candle points the same way as the published imbalance than on days when it points the
other way. A pre-registered study (plan locked 2026-11-10 19:37 UTC, one registered run) tested this prediction on the rule's
644 signal days of 2021--2026, the days of its holdout.

For each year the study used the 50 largest common-stock positions of SPY at the previous
year-end, from the fund's N-PORT reports to the SEC; they held between 51\% and 63\% of its net assets. On each signal day it
bought from Databento the imbalance messages that each stock's listing exchange sent between 15:50 and 16:00 (NYSE
Pillar for NYSE-listed stocks, Nasdaq TotalView-ITCH for Nasdaq-listed ones). For each stock it kept the first
closing-auction imbalance record received in the 15:50 minute and summed the signed dollar imbalances,
$D = \sum_i s_i Q_i p_i$, where $s_i$ is $+1$ for a buy imbalance, $-1$ for a sell imbalance and 0 for none, $Q_i$
is the imbalance in shares and $p_i$ the reference price, in exact integer arithmetic. A day with at least 10 usable
stocks and $D \ne 0$ is classified, and it \emph{agrees} if the sign of $D$ is the candle's direction. The outcome
is the rule's own: a hit when the CFD moves from the 15:51 open to the 15:59 close in the candle's direction. The
primary test (H1) compares the hit rates on agreeing and disagreeing days, $d = h_a - h_d$, and passes if $d > 0$
and $t = d/\mathrm{se} \ge 2.0$, with the standard error of two independent proportions. Two secondary tests,
which decide nothing, ask whether the candle agrees with the imbalance on more than half of the days (H2) and
whether the imbalance direction alone predicts the move (H3), with a Holm correction. A block bootstrap gives an
interval for $d$, and the robustness checks of Table~\ref{tab:imbtest} were fixed in the plan, as was the wording
of the verdict. By the plan's power calculation, which treats days as independent, 80\% power at the bar needs a
true $d$ of about 11 to 13 points, so the test can detect only a large
concentration of the edge on agreeing days.

Every signal day was classified. All 50 stocks had a usable first record on each of the 644 days, and the first
closing record of the minute arrived a median of about 0.1 seconds after 15:50:00 on both exchanges (at most
1.1 seconds). 

\begin{table}[htbp]
\centering
\caption{The closing-imbalance test: the rule's hit rate on days that agree and disagree with the 15:50 imbalance.}
\label{tab:imbtest}
\begin{tabular}{@{}lrrrr@{}}
\toprule
 & Agreeing days & Disagreeing days & Difference & Statistic \\
\midrule
\multicolumn{5}{@{}l}{\textit{Panel A: registered tests}} \\
H1: hit rate (primary) & 55.9\% (469) & 61.7\% (175) & $-5.9$ & $t = -1.35$ \\
\quad Block-bootstrap 95\% interval & & & \multicolumn{2}{l}{$-14.8$ to $+4.0$} \\
H2: candle agrees with imbalance & \multicolumn{2}{c}{469 of 644 (72.8\%)} & & $p = 3.1\times10^{-32}$ \\
H3: imbalance direction right & \multicolumn{2}{c}{329 of 644 (51.1\%)} & & $p = 0.30$ \\
\midrule
\multicolumn{5}{@{}l}{\textit{Panel B: robustness}} \\
Latest record before 15:51 & 55.6\% (459) & 62.2\% (185) & $-6.6$ & $t = -1.55$ \\
NYSE-listed stocks only & 56.2\% (372) & 59.2\% (272) & $-3.0$ & $t = -0.76$ \\
Nasdaq-listed stocks only & 57.3\% (499) & 57.9\% (145) & $-0.6$ & $t = -0.13$ \\
Adjusted for year and direction & & & $-6.4$ & 12 cells \\
Move in the candle's direction (bp) & $2.41$ (469) & $2.24$ (175) & $+0.16$ & $t = 0.13$ \\
Option P\&L, \$ a trade & $+\dollar80.43$ (423) & $+\dollar33.58$ (162) & $+\dollar46.85$ & $t = 0.79$ \\
\bottomrule
\end{tabular}
\tabnotes{Signal days 2021--2026; all 644 were classified. Cells give the value and, in brackets, the number of
days (trades for the option P\&L). Differences of hit rates are in percentage points; $t$ uses the standard error
of two independent proportions, and the move and option P\&L rows use Welch's $t$. The bootstrap draws 20-day blocks
of the 1{,}446 sessions, 10{,}000 of 10{,}000 draws kept. H2 and H3 are exact one-sided binomial tests
against one half (the half benchmark is descriptive), with a Holm correction: H2 rejects,
H3 does not reject. Panel~B was fixed in the plan and cannot change the verdict.
Move: from the 15:51 open to the 15:59 close in the candle's direction, basis points. Option P\&L: \ITM{10} at the
15:51 ask, held to settlement, \$0.09 fee, on the 585 priced trades.}
\end{table}

\begin{table}[htbp]
\centering
\caption{Signal days that agree and disagree with the 15:50 closing imbalance, by year.}
\label{tab:imbyear}
\begin{tabular}{@{}lrrlrrlr@{}}
\toprule
 & & \multicolumn{3}{c}{Agreeing days} & \multicolumn{3}{c}{Disagreeing days} \\
\cmidrule(lr){3-5}\cmidrule(l){6-8}
Year & Days & $n$ & Hit rate (95\% CI) & bp & $n$ & Hit rate (95\% CI) & bp \\
\midrule
2021 & 111 & 92 & 62.0\% (52--71) & $4.67$ & 19 & 78.9\% (57--91) & $5.12$ \\
2022 & 105 & 76 & 57.9\% (47--68) & $4.00$ & 29 & 69.0\% (51--83) & $3.88$ \\
2023 & 113 & 73 & 53.4\% (42--64) & $0.33$ & 40 & 55.0\% (40--69) & $1.18$ \\
2024 & 106 & 69 & 59.4\% (48--70) & $2.59$ & 37 & 43.2\% (29--59) & $-1.90$ \\
2025 & 117 & 91 & 53.8\% (44--64) & $1.39$ & 26 & 69.2\% (50--83) & $3.97$ \\
2026 & 92 & 68 & 47.1\% (36--59) & $0.96$ & 24 & 70.8\% (51--85) & $4.29$ \\
\midrule
All & 644 & 469 & 55.9\% (51--60) & $2.41$ & 175 & 61.7\% (54--69) & $2.24$ \\
\bottomrule
\end{tabular}
\tabnotes{Signal days of the rule (a 15:50 body of 60\% or more), 2021--2026. A day agrees when the sign of the
summed dollar imbalance of the 50 largest stocks, from each stock's first closing-imbalance record received in the
15:50 minute, equals the candle's direction. Every signal day was classified, so no day is unclassified. Hit rate:
share of days on which the CFD moved from the 15:51 open to the 15:59 close in the candle's direction, with its
Wilson 95\% interval in percent; bp: the mean of that move in the candle's direction, in basis points. Over all
days the difference in hit rates is $-5.9$ points ($t = -1.35$), with a 20-day block-bootstrap 95\%
interval from $-14.8$ to $+4.0$ points.}
\end{table}

Table~\ref{tab:imbtest} reports the results. The candle and the published
imbalance usually pointed the same way: they agreed on 469 of the 644 days
(72.8\%), far more often than half (H2, $p = 3.1\times10^{-32}$), but agreement did not
make the candle more accurate. It was right on 55.9\% of
agreeing days and on 61.7\% of disagreeing days, so $d = -5.9$ points (se
4.3, $t = -1.35$), and H1 did not pass. The block-bootstrap interval of $d$, from
$-14.8$ to $+4.0$ points, lies below the 11 to 13 points that the test was designed to detect.
The imbalance direction alone was right on 51.1\% of days (H3, $p = 0.30$). The robustness
checks agree: the difference in hit rates stayed negative with the latest record of the minute, with NYSE or Nasdaq
stocks alone and within year and direction. The average move in the candle's direction was about the same on both
kinds of day ($2.41$ against $2.24$ basis points), and the option trade made more on agreeing
days ($+\dollar80.43$ against $+\dollar33.58$ a trade), but not significantly ($t = 0.79$). By year
(Table~\ref{tab:imbyear}), disagreeing days had the higher hit rate in five of the six years, all but 2024.

In the words the plan fixed in advance, the prediction was not confirmed, the
test has limited power, and a failure does not refute the auction explanation. The price may adjust within the 15:50
candle itself and leave little after 15:51, imbalances change after their first publication, the stocks'
imbalances can offset one another in $D$, and 50 stocks are only part of the index. The high agreement is itself an
association, not a cause: the candle and the imbalance can both reflect the same closing orders, such as index
funds' orders known earlier in the day, and because records are timed by when Databento received them, the study
cannot show that the publication came before the price move. What it does show is narrower: on these days the
rule's edge was not concentrated on the days on which its candle agreed with the imbalance first published for the
50 largest stocks. Whether the effect changed when the NYSE moved its publication to 15:50 in 2019, and whether it
depends on the size of the imbalance, were outside the plan.

% =====================================================================
\section{Forward evidence}
\label{sec:forward}



\begin{table}[htbp]
\centering
\caption{Forward signals at real Robinhood prices, 22--30 October 2026.}
\label{tab:fwdstrikes}
\begin{tabular}{lccrr}
\toprule
Strike & trades & wins & total, hold to close & total, sell 15:59 \\
\midrule
\ITM{20} & 4 & 3 & $+\dollar1{,}493.64$ & $+\dollar1{,}329.28$ \\
\ITM{15} & 4 & 3 & $+\dollar1{,}193.64$ & $+\dollar1{,}084.28$ \\
\ITM{10} & 4 & 3 & $+\dollar1{,}213.64$ & $+\dollar1{,}164.28$ \\
\ITM{5} & 4 & 3 & $+\dollar658.64$ & $+\dollar484.28$ \\
\ATM & 4 & 2 & $-\dollar69.36$ & $-\dollar173.72$ \\
\OTM{5} & 4 & 1 & $-\dollar682.36$ & $-\dollar763.72$ \\
\OTM{10} & 4 & 0 & $-\dollar529.36$ & $-\dollar565.72$ \\
\OTM{15} & 4 & 0 & $-\dollar305.36$ & $-\dollar341.72$ \\
\bottomrule
\end{tabular}
\tabnotes{Strong-body signal, four trades, \$0.09 fee, dollars per contract. Four
trades support no conclusion.}
\end{table}

From 21 October to 6 November 2026, the strong-body candle fired on seven of
thirteen sessions and was right on six. Only four of those days have real
one-minute option prices; on the others the broker's minute bars had expired
before they were collected. Table~\ref{tab:fwdstrikes} prices the four trades at
every strike. Three of the four signals were right, by 5 to 8 points, which was
enough for in-the-money options but not for out-of-the-money ones, and selling at
15:59 was worse than holding at every strike. Four trades are too few to conclude anything, and at the recent effect size,
confirming the option P\&L would take many years of
forward signals.

% =====================================================================
\section{Discussion}
\label{sec:discussion}



\subsection{What the evidence supports}

After all of these studies the trading rule itself is unchanged, although the
evidence behind it has changed. The rule, as tested, is to check at 15:51 whether the 15:50
one-minute candle has a body of at least 60\% of its range and, if it does, to buy
one SPXW 0DTE option about 10 points in the money in the candle's direction (a
call 10 points below the rounded index after an up candle, a put 10 points above
it after a down candle) and hold it to settlement. The candle was always read
from the CFD feed. The signal passed a pre-registered holdout on index moves, and
at real asks the rule made money on 2021--2026, about two-thirds as much as the
model implied. Two searches over 2{,}178 nearby settings found nothing reliably
better, and a third over 7{,}128 settings that included other closing candles
chose the 15:50 candle again. The re-validation found that the effect begins near
a 60\% body but cannot rank the thresholds between 60\% and about 75\%; beyond
the threshold, a larger ratio was not followed by a larger gain. Deep
in-the-money strikes were the most consistent, in the model and at real prices,
and no early exit improved on holding.

The edge is small, about two index points over nine minutes, or roughly six
cents per dollar of premium. It is concentrated in 2021--2022 and is not
significant on its own since 2023. It is also fragile to execution: the bid-ask
spread erased it before 2021, a one-minute delay costs a meaningful share of
it, and on smaller contracts the spread exceeds it. Confirming the current rule
does not show that it is the best rule, and two failed searches for something better do not show that
nothing better exists.

The confirmation also applies only to the CFD candle. Read from the SPX index
candle, which is what a trader sees, the rule made $+\$39.44$ a trade against
$+\$67.24$ on the same days and did not meet its pre-registered
criterion (Section~\ref{sec:realindex}). The index candle fires about as often as
the tested rule only at a threshold of roughly 70\%, and even above 70\% the days
on which it fired alone lost money. This is a practical problem as well as a statistical one, because in this
project the CFD candle comes from daily history files
that appear after the session, so the live alert could at first read only the
index candle at 15:51.

The SPY candle partly solves that problem. Read from it, the rule made $+\$71.74$ a trade,
about what the CFD candle made on the same days, and met its retrospective
criterion (Section~\ref{sec:spycandle}); since that study the live alert reads
the SPY candle too and acts on a rule built from it, the SPY candle at 70\% or
more with the index candle at 60\% or more in the same direction. Because SPY
was the third source tried on the same outcomes, and one of its lines was chosen after agreement between candles had been seen to
do well, I treat this as a lead for the forward record and not as a confirmation. Across
the three sources the pattern is the same. The gain came on sessions on which
every candle showed a strong body, and a strong body on only some of them
carried no edge.

The search of Section~\ref{sec:search} is a comparison and not a test. No
other closing minute, cut, direction or strike beat the 15:50 candle on
2021--2023, and the settings that did best over 2021--2026 were all 15:50
candles traded in their direction, most with a deeper strike than \ITM{10}. Its
own choice was not clearly beyond luck, and because it fell at 15:50 it could not
be tested on the later years. Those settings are hypotheses for the forward
record.

The test of Section~\ref{sec:imbalance} is about why the signal exists, not whether it pays. The candle usually points the same way as the closing-auction
imbalance published at 15:50, but it was not more accurate on the days on which it
did, so the most direct prediction of the auction explanation was not confirmed.
The test could detect only a large difference, so that explanation is still open.

\subsection{Negative results}
\label{sec:fails}

This subsection lists the failures in one place. Of the 41 candle features, 38
did not qualify, and one of the three that did, a whole-day feature, failed out
of sample. The
plain 15:50 direction, without a body filter, predicts the close but not
accurately enough to pay for an option. The pre-registered skew rule failed for
far out-of-the-money strikes. No early exit, momentum filter or dealer-gamma
condition improved the trade, and all 26 real-time exit rules tested at real
bids earned less than holding. At real asks the trade lost money before 2021.
The first two grid searches failed to find a better setting, the second because
its winner failed the luck diagnostic, and the third chose a 15:50 setting whose
luck diagnostic missed its bar. A limit order at the mid did not capture the
spread, and the trade was not confirmed on XSP or SPY. Read from the index candle,
the rule did not meet its criterion, and its 60--70\% band showed no edge, as on
the SPY candle. The candle was not more accurate on days on which it agreed with
the published closing-auction imbalance. Finally, the forward test lost three
signal days because the broker's minute bars expired before they were
collected.

\subsection{Limitations}
\label{sec:limits}

The design has several limitations, and the first is that there was no blind
test. The passing candle overlaps a known effect, and the 2021--2026 index
outcomes were reused by the real-price confirmation and by the later descriptive
analyses.

The signal also comes from a CFD. Every trading decision uses a Dukascopy CFD candle and not the SPX index
itself, and only the option prices and the settlement are real. I measured how
much this matters. Read from
the index candle, the rule made $+\$39.44$ a trade and did not meet its
criterion, the 60--70\% band showed no edge, and the days on which only the index
candle fired lost money (Section~\ref{sec:realindex}). That check reused outcomes
that had already been seen. Read from the SPY candle, the rule met a
retrospective criterion, but SPY was the third source tried on those outcomes,
and a second vendor's SPY candle differed by a cent or two on about a third of
the sessions compared (Section~\ref{sec:spycandle}).

Execution adds further uncertainty. The real-price studies assume a fill at the
quoted ask at the open of the minute, which ignores order size and fast markets.
Trade bars show that the contract traded near the ask, but not the size or the
queue, and the limit-order results are proxies that assume no delay. Entry is
assumed at 15:51:00, while a person reacting to the candle trades seconds later.
The trade also carries risk: a deep in-the-money contract moves almost point for
point with the index, and at real asks the worst single trade on 2021--2026 lost
\$2{,}210 at \ITM{10}.

The statistical tools and the samples have limits of their own. The luck
diagnostic is approximate. Expiry frequency, volatility and pricing changed over
the samples, so a circular shift is only an approximate benchmark for ``no
information'', yet it decided the re-tuning study and was one of the gates of the
final search. The early sample is thin: before 2016 SPXW had a same-day expiry
only on some days, and daily expiries began in 2022, so tuning on 2013--2020
rests on a thin, uneven sample, and the 2021--2023 re-tuning on only three years.
The strike and exit studies use a model calibrated on seven days of 2026 data,
and the exit rules tested at real bids (Section~\ref{sec:rtexit}) use one-minute
snapshots and assume one contract and a firm bid one minute after the decision.
Finally, the forward sample is small, with four priced forward trades.

% =====================================================================
\section{Conclusion}
\label{sec:conclusions}

I asked whether a one-minute candle near the close predicts the
S\&P~500 into the close, and whether the prediction can be traded with same-day
options once the spread is paid. One candle passed a pre-registered, out-of-sample test, the 15:50 one-minute
SPX candle with a strong body, traded in its own direction to the close. It is
a stricter version of an effect that was already known, and it has faded since
2023. Traded with a deep in-the-money 0DTE option bought at the quoted ask and
held to settlement, it made money on 2021--2026, though less than a pricing
model implied. The largest cost is the bid-ask spread, which took the whole
edge before 2021, and a limit order at the mid does not recover it because it
misses the best days. Re-tuned settings, early exits, momentum filters and
smaller contracts did not improve on it, and a search over other closing candles chose the 15:50 candle
again.

The evidence has limits. The edge is small, most of it comes
from 2021--2022, and it is fragile to execution. The confirmation holds for the
candle as the backtests read it, from a CFD that tracks the index. Read from the
SPX index candle itself, under a plan locked in advance, the rule earned less and
did not meet its criterion; the gain was confined to days on which both candles
showed a strong body. Read from the SPY candle, which a trader can watch live, it
met a retrospective criterion, with the gain again confined to sessions on which
all three candles showed a strong body. A pre-registered test with the
exchanges' imbalance messages found that the candle usually agrees with the
closing-auction imbalance published at 15:50 but is no more accurate when it
does, so that explanation of the signal was not confirmed.

One advantage of pre-registering every study is that the failed tests can be
reported and used. Because each plan was fixed and hashed before the data it
governs was read, a reader can see which refinements were tried and did not
help, and which explanation the data did not confirm. To summarise, the signal
passed a test at real prices when read from the CFD candle, did not pass when
read from the index candle, and met a retrospective criterion when read from
the SPY candle. It remains sensitive to execution. I think it is worth
continuing to measure, but it is not an established edge.

There are three things I would do next. Settings suggested by the tuning results,
such as the 15:48 candle, and the stricter 15:50 settings favoured by the final
search should be treated as hypotheses and tested only on days after November
2026, under new registered plans. The forward record begun in November 2026
(Section~\ref{sec:audit}), which since the SPY study holds all three candles and
tracks the rule the alert now acts on, should run until it can test, on new days,
whether sessions on which all three candles show a strong body continue to beat
sessions on which only some do. The SPY candle's retrospective pass is a reason to keep that record and does
not replace it, and the 80\% cut on the index candle and the futures-lead
explanation of Section~\ref{sec:realindex} are hypotheses for it. Finally,
forward collection should continue, with the number of forward days needed for
a test fixed in advance, ideally alongside records of actual order fills,
because only real fills can show how the trade should be executed.

\paragraph{Data and code availability.} The analysis code, the pre-registration
plans, their hash ledger and the shorter journal version of this paper are public
at \url{https://github.com/kv1514/spx-closing-candle-study}. The OPRA data are licensed
from Databento and are not redistributed.

% =====================================================================
\begingroup
\setstretch{1.0}
\setlength{\bibsep}{6pt}
\begin{thebibliography}{99}

\bibitem[Bachelier(1900)]{bachelier1900}
Bachelier, L. (1900). Th\'eorie de la sp\'eculation. \textit{Annales
scientifiques de l'\'Ecole Normale Sup\'erieure}, 17:21--86.

\bibitem[Baltussen et al.(2021)]{bdlm2021}
Baltussen, G., Z. Da, S. Lammers, and M. Martens (2021). Hedging demand and
market intraday momentum. \textit{Journal of Financial Economics}, 142:377--403.

\bibitem[Gao et al.(2018)]{ghlz2018}
Gao, L., Y. Han, S. Z. Li, and G. Zhou (2018). Market intraday momentum.
\textit{Journal of Financial Economics}, 129:394--414.

\bibitem[Harvey et al.(2016)]{harvey2016}
Harvey, C. R., Y. Liu, and H. Zhu (2016). \ldots\ and the cross-section of
expected returns. \textit{Review of Financial Studies}, 29:5--68.

\bibitem[Holm(1979)]{holm1979}
Holm, S. (1979). A simple sequentially rejective multiple test procedure.
\textit{Scandinavian Journal of Statistics}, 6:65--70.

\bibitem[Newey and West(1987)]{neweywest1987}
Newey, W. K. and K. D. West (1987). A simple, positive semi-definite,
heteroskedasticity and autocorrelation consistent covariance matrix.
\textit{Econometrica}, 55:703--708.

\bibitem[Nosek et al.(2018)]{nosek2018}
Nosek, B. A., C. R. Ebersole, A. C. DeHaven, and D. T. Mellor (2018). The
preregistration revolution. \textit{Proceedings of the National Academy of
Sciences}, 115:2600--2606.

\bibitem[Wilson(1927)]{wilson1927}
Wilson, E. B. (1927). Probable inference, the law of succession, and statistical
inference. \textit{Journal of the American Statistical Association}, 22:209--212.

\end{thebibliography}
\endgroup

% =====================================================================
\appendix
\setlength{\intextsep}{6pt plus 3pt minus 2pt}
\counterwithin{table}{section}
\renewcommand{\thetable}{\thesection\arabic{table}}
\titleformat{\section}{\normalfont\Large\bfseries}{Appendix \thesection:}{0.5em}{}

\section{Pricing model calibration}
\label{app:cal}

This appendix sets out how the Bachelier model of Section~\ref{sec:pricing} was
calibrated and reports the model variants used in Section~\ref{sec:model}.

\subsection{Volatility and skew}

For a day with one-minute realised volatility $RV$, the root mean square of the
09:31--15:50 log returns, the model sets
\begin{equation}
\sigma' = k_{D1}\cdot RV \cdot 3 \cdot L,
\end{equation}
where $3 = \sqrt{9}$ covers the nine minutes from 15:51 to 15:59 and
$L = 7{,}709.58$ is the common index level. A strike's price is $f_j \cdot C$ or
$f_j \cdot P$, where $f_j$ is the skew factor for offset $j$. The rules came
first and the values from the data. Implied volatility was fitted by least squares
to the 15:51 opens of the out-of-the-money calls and puts within 30 points, on the
seven days with real one-minute option bars (2026-10-20 to 10-30). $k_{RH}$ is the
median ratio of implied to realised volatility on Robinhood's SPX bars, $c$ is the
median ratio of Robinhood to CFD realised volatility over 22 overlapping days,
and $k_{D1} = k_{RH}\,c$. The skew factor $f_j$ is the median ratio of market to
model price across days and sides (Table~\ref{tab:cal}).

\subsection{Model variants}

Variant (a) is this calibration. Variant (b) uses a single constant volatility,
variant (c) the scan's own cost assumption, and variant (d) calibrates $k$ on
10-minute bars. The sensitivities change the half-spread, scale $k$, or drop the
skew factors (Table~\ref{tab:variants}). ``Credible'' variants in the text are all
but (b), which ignores how volatile each day was.

\begin{table}[htbp]
\centering
\caption{Calibrated skew factors.}
\label{tab:cal}
\begin{tabular}{lrrrrrrrrr}
\toprule
& \ITM{20} & \ITM{15} & \ITM{10} & \ITM{5} & \ATM & \OTM{5} & \OTM{10} & \OTM{15} & \OTM{20} \\
\midrule
$f_o$ & 1.005 & 1.020 & 1.005 & 1.007 & 0.969 & 1.071 & 2.091 & 2.256 & 54.000 \\
ratios & 10 & 13 & 13 & 14 & 14 & 14 & 13 & 13 & 10 \\
\bottomrule
\end{tabular}
\tabnotes{``Ratios'' is the number of market-to-model price ratios behind each
factor. $k_{RH} = 1.089$; $c = 0.978$ over 22 days; $k_{D1} = 1.065$; ten-minute
$k = 1.319$ over 22 days. The variant (a) model move to the close has a median of
10.10 bps (10th--90th percentile 5.72--19.20). Variant (c) uses
$\sigma = 10.78$ points.}
\end{table}

\begin{table}[htbp]
\centering
\caption{Mean modelled P\&L per contract in every variant.}
\label{tab:variants}
\begin{tabular}{lrrrrrrr}
\toprule
Variant & \ITM{20} & \ITM{15} & \ITM{10} & \ITM{5} & \ATM & \OTM{5} & \OTM{10} \\
\midrule
(a) calibrated & $+162$ & $+143$ & $+166$ & $+161$ & $+155$ & $+92$ & $-34$ \\
(b) constant vol & $+192$ & $+192$ & $+242$ & $+262$ & $+261$ & $+201$ & $+127$ \\
(c) realized-move-fair & $+186$ & $+181$ & $+158$ & $+122$ & $+83$ & $+64$ & $+62$ \\
(d) ten-minute $k$ & $+136$ & $+104$ & $+110$ & $+85$ & $+72$ & $+12$ & $-149$ \\
(a) half-spread \dollar0.05 & $+167$ & $+148$ & $+171$ & $+166$ & $+160$ & $+97$ & $-29$ \\
(a) half-spread \dollar0.20 & $+152$ & $+133$ & $+156$ & $+151$ & $+145$ & $+82$ & $-44$ \\
(a) $k \times 0.75$ & $+180$ & $+172$ & $+214$ & $+235$ & $+239$ & $+170$ & $+64$ \\
(a) $k \times 1.25$ & $+135$ & $+102$ & $+107$ & $+81$ & $+68$ & $+8$ & $-155$ \\
(a) no skew & $+172$ & $+174$ & $+171$ & $+166$ & $+143$ & $+106$ & $+72$ \\
\bottomrule
\end{tabular}
\tabnotes{Level-normalised dollars, \$0.09 fee, 644 modelled trades. Strikes 10--20
points in the money stay between $+\$102$ and $+\$242$; at the money ranges from
$+\$68$ to $+\$261$.}
\end{table}

\needspace{10\baselineskip}
\section{Additional tables}
\label{app:tables}

This appendix reports three tables that support Sections~\ref{sec:refine}
and~\ref{sec:momentum}.

\subsection{The body threshold profile}

Table~\ref{tab:threshold} gives the profile behind Figure~\ref{fig:threshold}
in Section~\ref{sec:refine}.

\begin{table}[htbp]
\centering
\renewcommand{\arraystretch}{0.92}
\caption{The body threshold profile.}
\label{tab:threshold}
\setlength{\tabcolsep}{4pt}
\begin{tabular}{lrrrrrrrr}
\toprule
& \multicolumn{4}{c}{2021--2023 (tuning)} & \multicolumn{4}{c}{2024--2026 (after the decision)} \\
\cmidrule(lr){2-5}\cmidrule(lr){6-9}
Body at least & trades & \$/day & \$/trade & won & trades & \$/day & \$/trade & won \\
\midrule
0.40 & 407 & $+7.23$ & $+11.09$ & 53\% & 468 & $+34.63$ & $+51.57$ & 51\% \\
0.45 & 376 & $+12.60$ & $+20.92$ & 53\% & 437 & $+33.24$ & $+53.02$ & 51\% \\
0.50 & 341 & $+9.08$ & $+16.61$ & 52\% & 393 & $+35.99$ & $+63.82$ & 52\% \\
0.55 & 303 & $+15.60$ & $+32.12$ & 53\% & 353 & $+30.04$ & $+59.32$ & 52\% \\
\textbf{0.60 (current)} & \textbf{270} & $\mathbf{+29.76}$ & $\mathbf{+68.78}$ & \textbf{55\%} & \textbf{315} & $\mathbf{+29.97}$ & $\mathbf{+66.32}$ & \textbf{53\%} \\
0.65 & 227 & $+26.38$ & $+72.62$ & 55\% & 281 & $+26.27$ & $+65.15$ & 53\% \\
0.70 & 197 & $+25.90$ & $+82.16$ & 57\% & 238 & $+32.16$ & $+94.20$ & 55\% \\
0.75 & 173 & $+28.43$ & $+102.72$ & 58\% & 207 & $+24.56$ & $+82.71$ & 53\% \\
0.80 & 133 & $+37.84$ & $+177.82$ & 62\% & 167 & $+13.26$ & $+55.33$ & 50\% \\
0.85 & 89 & $+25.51$ & $+179.17$ & 63\% & 130 & $+8.35$ & $+44.77$ & 48\% \\
0.90 & 56 & $+11.19$ & $+124.89$ & 62\% & 95 & $+16.36$ & $+120.02$ & 55\% \\
\bottomrule
\end{tabular}
\tabnotes{15:50 candle, with the candle, \ITM{10}, real asks, \$0.09 fee, dollars
per contract. Descriptive. The 2024--2026 columns were examined only after the
re-tuning decision was recorded. At 0.85 and 0.90 the setting could not be
eligible (fewer than 100 tuning trades).}
\end{table}

\subsection{Momentum filters}

Table~\ref{tab:momB} reports the four momentum filters of Part B of the study in
Section~\ref{sec:momentum}.

\begin{table}[htbp]
\centering
\renewcommand{\arraystretch}{0.92}
\caption{Momentum filters on the candle trade (momentum study, Part B).}
\label{tab:momB}
\setlength{\tabcolsep}{4pt}
\begin{tabular}{lrrrrrrrr}
\toprule
& \multicolumn{2}{c}{trades} & \multicolumn{2}{c}{mean P\&L / trade} & \multicolumn{4}{c}{mean difference from R0 per day} \\
\cmidrule(lr){2-3}\cmidrule(lr){4-5}\cmidrule(lr){6-9}
Rule & kept & skipped & kept & skipped & \$/day ($t$) & 2013--16 & 2017--20 & index (bps) \\
\midrule
\multicolumn{9}{l}{\emph{2013--2020: discovery, the selection test}} \\
R0 (all) & 905 & 0 & $+133.21$ & --- & --- & --- & --- & --- \\
R1 & 490 & 415 & $+198.58$ & $+56.02$ & $-12.65$ ($-1.39$) & $+1.27$ & $-24.45$ & $-0.320$ \\
R2 & 471 & 434 & $+202.18$ & $+58.36$ & $-13.78$ ($-1.44$) & $+1.39$ & $-26.64$ & $-0.327$ \\
R3 & 376 & 529 & $+252.84$ & $+48.18$ & $-13.87$ ($-1.37$) & $+1.05$ & $-26.51$ & $-0.378$ \\
R4 & 707 & 198 & $+149.78$ & $+74.04$ & $-7.98$ ($-1.06$) & $+5.86$ & $-19.70$ & $-0.200$ \\
\midrule
\multicolumn{9}{l}{\emph{2021--2026: descriptive, not a test}} \\
R0 (all) & 644 & 0 & $+166.08$ & --- & --- & & & --- \\
R1 & 312 & 332 & $+217.22$ & $+118.02$ & $-27.10$ ($-2.09$) & & & $-0.383$ \\
R2 & 315 & 329 & $+184.64$ & $+148.31$ & $-33.74$ ($-2.63$) & & & $-0.491$ \\
R3 & 225 & 419 & $+210.59$ & $+142.18$ & $-41.20$ ($-2.80$) & & & $-0.579$ \\
R4 & 458 & 186 & $+190.44$ & $+106.09$ & $-13.65$ ($-1.37$) & & & $-0.213$ \\
\bottomrule
\end{tabular}
\tabnotes{\ITM{10}, model variant (a), \$0.09 fee, level-normalised dollars per
contract. R0 takes every trade; R1 takes it only when it agrees with
$r_{\text{rod}}$, R2 with the first half-hour return, R3 with both; R4 skips it on
strong disagreement with $r_{\text{rod}}$. 2013--2020 uses 1{,}838 paired days
and 2021--2026 uses 1{,}446. No rule was eligible, so no confirmation test was
run.}
\end{table}

\needspace{14\baselineskip}
\subsection{Dealer gamma}

Table~\ref{tab:partc} reports the two dealer-gamma tests of Part C of the same
study.

\begin{table}[htbp]
\centering
\renewcommand{\arraystretch}{0.92}
\caption{Dealer gamma (momentum study, Part C).}
\label{tab:partc}
\begin{tabular}{@{}llrrrl@{}}
\toprule
Test & Sample & Coefficient & $t$ & $p$ & Verdict \\
\midrule
C1, $d$ & 2013--2026, 3{,}238 days & $+0.022$ & 0.95 & 0.34 & not supported \\
C2, $g$ & 2013--2026, 1{,}549 trades & $+\dollar169$ & 2.18 & 0.029 & not supported \\
C2 guard (index, bps) & same trades & $+3.13$ & 2.54 & 0.011 & --- \\
\midrule
C2, $g$ & 2013--2020, 905 trades & $+\dollar240$ & 1.83 & & descriptive \\
C2, $g$ & 2021--2026, 644 trades & $+\dollar84$ & 0.96 & & descriptive \\
Real asks, difference & 2021--2026, 585 trades & $-\dollar16$ & $-0.25$\textsuperscript{a} & & descriptive \\
\bottomrule
\end{tabular}
\tabnotes{C1: $r_{\text{last}}$ on $r_{\text{rod}}$, the negative-gamma indicator
$N$ and $N\cdot r_{\text{rod}}$. C2: the candle trade's model P\&L (\ITM{10}, \$0.09
fee) on $N$. Newey--West $t$; the proxy was negative on 30\% of days. The real-ask
row is descriptive and not in the plan; \textsuperscript{a}two-sample $t$, trades
independent.}
\end{table}

\end{document}
